\documentclass{beamer}
\usetheme{Madrid}
\usecolortheme{default}
\usepackage[utf8]{inputenc}
\usepackage{amsmath}
\usepackage{amsfonts}
\usepackage{amssymb}
\usepackage{tikz}

\title{CRISC Exam - Hard Questions with Explanations}
\subtitle{From Open Source CRISC Exam Guide}
\author{Compiled from GTPs and AI}
\date{\today}

\begin{document}
	
	\begin{frame}
		\titlepage
	\end{frame}
	
	\begin{frame}{Outline}
		\tableofcontents
	\end{frame}
	
	% ================ SECTION 1: DOMAIN 1 - GOVERNANCE ================
	
	\section{Domain 1: Governance}
	
\begin{frame}{Q1: Governance Accountability}
	\begin{block}{Question}
		Which is the PRIMARY objective of an organizational governance structure?
	\end{block}
	
	\begin{enumerate}[(A)]
		\item Monitor the effectiveness of operational controls
		\item Establish accountability and direction for achieving enterprise objectives
		\item Manage business processes and operational risks
		\item Implement risk response activities across the enterprise
	\end{enumerate}
\end{frame}

\begin{frame}{Q1: Explanation}
	\begin{block}{Correct Answer: B}
		Governance provides direction, oversight, and accountability
		to ensure enterprise objectives are achieved and aligned with
		stakeholder expectations.
	\end{block}
	
	\begin{block}{Option Analysis}
		\begin{itemize}
			\item \textbf{A - Incorrect:} Monitoring controls is primarily a management function.
			\item \textbf{B - Correct:} Governance establishes accountability, oversight, and strategic alignment.
			\item \textbf{C - Incorrect:} Managing processes and risks is an operational responsibility.
			\item \textbf{D - Incorrect:} Risk response implementation is performed by management.
		\end{itemize}
	\end{block}
\end{frame}


\begin{frame}{Q1: Governance Versus Management} \begin{block}{Question} Senior leadership wants to ensure risk management activities support organizational objectives and stakeholder expectations. Which is MOST important? \end{block} \begin{enumerate}[(A)] \item Defining operational procedures \item Establishing governance and accountability structures \item Performing periodic risk assessments \item Implementing compensating controls \end{enumerate} \end{frame} \begin{frame}{Q1: Correct Answer} \begin{block}{Answer: B} Governance provides direction, oversight, and accountability to ensure that risk management activities support business objectives and stakeholder expectations. \end{block} \end{frame} \begin{frame}{Q1: Option Analysis} \textbf{A - Defining Operational Procedures (Incorrect)} \begin{itemize} \item Operational procedures describe how work is performed. \item This is a management responsibility. \item Procedures help execute strategy but do not establish enterprise accountability. \end{itemize} \vspace{0.2cm} \textbf{B - Establishing Governance and Accountability Structures (Correct)} \begin{itemize} \item Governance aligns risk management with business objectives. \item Establishes accountability and oversight. \item Ensures stakeholder expectations are addressed. \item Provides direction for management activities. \end{itemize} \end{frame} \begin{frame}{Q1: Option Analysis (Continued)} \textbf{C - Performing Periodic Risk Assessments (Incorrect)} \begin{itemize} \item Risk assessments identify and evaluate risks. \item They are important management activities. \item However, governance must be established before risk assessments can effectively support enterprise objectives. \end{itemize} \vspace{0.2cm} \textbf{D - Implementing Compensating Controls (Incorrect)} \begin{itemize} \item Compensating controls reduce risk when primary controls are unavailable or insufficient. \item Control implementation is a management responsibility. \item Controls support governance but do not replace it. \end{itemize} \end{frame} \begin{frame}{Q1: Exam Trap} \begin{alertblock}{Common CRISC Trap} Risk assessments and controls sound risk-focused and are often selected by candidates. However, the question emphasizes: \begin{itemize} \item Senior leadership \item Organizational objectives \item Stakeholder expectations \end{itemize} These are governance concepts, making governance and accountability structures the BEST answer. \end{alertblock} \end{frame} \begin{frame}{Q1: Governance vs Management} \begin{block}{Remember} \textbf{Governance} \begin{itemize} \item Provides direction \item Establishes accountability \item Aligns with stakeholder expectations \item Oversees performance \end{itemize} \textbf{Management} \begin{itemize} \item Implements controls \item Performs risk assessments \item Executes procedures \item Runs day-to-day operations \end{itemize} \end{block} \end{frame} \begin{frame}{Q1: Summary Question} \begin{block}{Remember for the Exam} Which activity provides the foundation that ensures risk management supports enterprise objectives? \end{block} \begin{enumerate} \item Operational Procedures \item Governance and Accountability Structures \textbf{(Correct)} \item Risk Assessments \item Compensating Controls \end{enumerate} \alert{Governance comes before management activities.} \end{frame}

	
	\begin{frame}{Q2: Risk Management First Step}
		\begin{block}{Question}
			Which of the following is the first step in IT risk management?
		\end{block}
		\begin{enumerate}[(A)]
			\item IT risk assessment
			\item Risk response and mitigation
			\item IT risk identification
			\item Risk reporting
		\end{enumerate}
	\end{frame}
	
	\begin{frame}{Q2: Explanation}
		\begin{block}{Correct Answer: C}
			IT risk identification is the first step of IT risk management. You cannot protect what you do not know exists.
		\end{block}
		\begin{block}{Option Analysis}
			\begin{itemize}
				\item \textbf{A - Incorrect:} Assessment comes after identification
				\item \textbf{B - Incorrect:} Response comes after assessment
				\item \textbf{C - Correct:} Identification is the first step
				\item \textbf{D - Incorrect:} Reporting is the final step
			\end{itemize}
		\end{block}
	\end{frame}
	
	\begin{frame}{Q3: Oversight Responsibility}
		\begin{block}{Question}
			Oversight and direction are the responsibility of who?
		\end{block}
		\begin{enumerate}[(A)]
			\item Management
			\item Senior staff
			\item IT risk manager
			\item Board of directors
		\end{enumerate}
	\end{frame}
	
	\begin{frame}{Q3: Explanation}
		\begin{block}{Correct Answer: D}
			Oversight and direction are the responsibility of the board of directors.
		\end{block}
		\begin{block}{Option Analysis}
			\begin{itemize}
				\item \textbf{A - Incorrect:} Management executes, doesn't provide oversight
				\item \textbf{B - Incorrect:} Senior staff don't provide governance oversight
				\item \textbf{C - Incorrect:} Risk manager executes risk activities
				\item \textbf{D - Correct:} Board provides oversight and direction
			\end{itemize}
		\end{block}
	\end{frame}
	
	\begin{frame}{Q4: Policy Documentation Hierarchy}
		\begin{block}{Question}
			The high-level intent of an organization's security practices is best established by which of the following?
		\end{block}
		\begin{enumerate}[(A)]
			\item Procedures
			\item Policies
			\item Guidelines
			\item Processes
		\end{enumerate}
	\end{frame}
	
	\begin{frame}{Q4: Explanation}
		\begin{block}{Correct Answer: B}
			Policies are the high-level intent of organizations for security practices.
		\end{block}
		\begin{block}{Option Analysis}
			\begin{itemize}
				\item \textbf{A - Incorrect:} Procedures are detailed steps
				\item \textbf{B - Correct:} Policies express management intent
				\item \textbf{C - Incorrect:} Guidelines are recommendations
				\item \textbf{D - Incorrect:} Processes are operational
			\end{itemize}
		\end{block}
	\end{frame}
	
	\begin{frame}{Q5: RACI Model}
		\begin{block}{Question}
			Which group should be reached out to for expert guidance on the working group?
		\end{block}
		\begin{enumerate}[(A)]
			\item Responsible
			\item Accountable
			\item Consulted
			\item Informed
		\end{enumerate}
	\end{frame}
	
	\begin{frame}{Q5: Explanation}
		\begin{block}{Correct Answer: C}
			Consulted working groups are the subject matter experts for guidance on initiatives in their domain of expertise.
		\end{block}
		\begin{block}{Option Analysis}
			\begin{itemize}
				\item \textbf{A - Incorrect:} Responsible does the work
				\item \textbf{B - Incorrect:} Accountable provides resources
				\item \textbf{C - Correct:} Consulted are subject matter experts
				\item \textbf{D - Incorrect:} Informed are kept updated
			\end{itemize}
		\end{block}
	\end{frame}
	
	\begin{frame}{Q6: Risk Culture Focus}
		\begin{block}{Question}
			An IT risk manager recently joined an organization with past security incidents. What should be the primary focus in the first few months?
		\end{block}
		\begin{enumerate}[(A)]
			\item Cite individuals who caused incidents
			\item Understand the organization's culture toward risk
			\item Draft strategy for IT risk management
			\item Complete all security awareness training
		\end{enumerate}
	\end{frame}
	
	\begin{frame}{Q6: Explanation}
		\begin{block}{Correct Answer: B}
			The risk manager needs to understand the culture of the organization before drafting strategy.
		\end{block}
		\begin{block}{Option Analysis}
			\begin{itemize}
				\item \textbf{A - Incorrect:} Blaming individuals is counterproductive
				\item \textbf{B - Correct:} Culture understanding is foundational
				\item \textbf{C - Incorrect:} Strategy comes after understanding culture
				\item \textbf{D - Incorrect:} Training is secondary to culture assessment
			\end{itemize}
		\end{block}
	\end{frame}
	
	
	
	\begin{frame}{Q8: Third LoD Responsibility}
		\begin{block}{Question}
			What is the primary responsibility of the third LoD?
		\end{block}
		\begin{enumerate}[(A)]
			\item Policy and procedure development
			\item Provide independent assurance of controls
			\item Perform periodic user reviews
			\item Restrict least privileged roles
		\end{enumerate}
	\end{frame}
	
	\begin{frame}{Q8: Explanation}
		\begin{block}{Correct Answer: B}
			The third LoD is primarily responsible for independent assurance of controls through internal/external audit.
		\end{block}
		\begin{block}{Option Analysis}
			\begin{itemize}
				\item \textbf{A - Incorrect:} Policy development is second LoD
				\item \textbf{B - Correct:} Independence is key for third LoD
				\item \textbf{C - Incorrect:} User reviews are operational (first LoD)
				\item \textbf{D - Incorrect:} Access control is operational
			\end{itemize}
		\end{block}
	\end{frame}
	

	
	
	\begin{frame}{Q11: Risk Tolerance vs Capacity}
		\begin{block}{Question}
			Which statement is correct?
		\end{block}
		\begin{enumerate}[(A)]
			\item Breaching risk tolerance threatens existence
			\item Breaching risk capacity threatens existence
			\item Risk tolerance and capacity are not related
			\item Risk tolerance and capacity are the same
		\end{enumerate}
	\end{frame}
	
	\begin{frame}{Q11: Explanation}
		\begin{block}{Correct Answer: B}
			Breaching risk capacity could threaten an organization's existence.
		\end{block}
		\begin{block}{Option Analysis}
			\begin{itemize}
				\item \textbf{A - Incorrect:} Tolerance breach is manageable
				\item \textbf{B - Correct:} Capacity breach threatens existence
				\item \textbf{C - Incorrect:} They are related
				\item \textbf{D - Incorrect:} They are different concepts
			\end{itemize}
		\end{block}
	\end{frame}
	
	\section{Domain 1: Legal and Ethics}
	
	\begin{frame}{Q12: HIPAA}
		\begin{block}{Question}
			Which federal law provides guidance on protecting sensitive health information?
		\end{block}
		\begin{enumerate}[(A)]
			\item CCPA
			\item HIPAA
			\item FFIEC
			\item GDPR
		\end{enumerate}
	\end{frame}
	
	\begin{frame}{Q12: Explanation}
		\begin{block}{Correct Answer: B}
			HIPAA is the federal law protecting sensitive health information.
		\end{block}
		\begin{block}{Option Analysis}
			\begin{itemize}
				\item \textbf{A - Incorrect:} CCPA is California privacy law
				\item \textbf{B - Correct:} HIPAA is health information law
				\item \textbf{C - Incorrect:} FFIEC is banking regulation
				\item \textbf{D - Incorrect:} GDPR is EU regulation
			\end{itemize}
		\end{block}
	\end{frame}
	
	\begin{frame}{Q13: CCPA Beneficiaries}
		\begin{block}{Question}
			Which group is a beneficiary of CCPA law?
		\end{block}
		\begin{enumerate}[(A)]
			\item US residents
			\item EU residents
			\item California residents
			\item Canada residents
		\end{enumerate}
	\end{frame}
	
	\begin{frame}{Q13: Explanation}
		\begin{block}{Correct Answer: C}
			CCPA protects California residents' personal information.
		\end{block}
		\begin{block}{Option Analysis}
			\begin{itemize}
				\item \textbf{A - Incorrect:} US residents is too broad
				\item \textbf{B - Incorrect:} EU residents are covered by GDPR
				\item \textbf{C - Correct:} California residents
				\item \textbf{D - Incorrect:} Canada residents have PIPEDA
			\end{itemize}
		\end{block}
	\end{frame}
	
	
	
	% ================ SECTION 2: DOMAIN 2 - IT RISK ASSESSMENT ================
	
	\section{Domain 2: Risk Management Life Cycle}

	\begin{frame}{Q17: Risk Identification Principle}
		\begin{block}{Question}
			Which risk management life cycle step emphasizes "You can't protect what you don't know exists"?
		\end{block}
		\begin{enumerate}[(A)]
			\item Risk and control monitoring
			\item Risk assessment
			\item Risk identification
			\item Risk categorization
		\end{enumerate}
	\end{frame}
	
	\begin{frame}{Q17: Explanation}
		\begin{block}{Correct Answer: C}
			Risk identification is the first step - you must know risks exist.
		\end{block}
		\begin{block}{Option Analysis}
			\begin{itemize}
				\item \textbf{A - Incorrect:} Monitoring happens after identification
				\item \textbf{B - Incorrect:} Assessment comes after identification
				\item \textbf{C - Correct:} Identification discovers risks
				\item \textbf{D - Incorrect:} Categorization comes after identification
			\end{itemize}
		\end{block}
	\end{frame}
	
	\begin{frame}{Q18: Formal Risk Assessment Requirement}
		\begin{block}{Question}
			Which is a formal requirement by many legal and regulatory compliance frameworks?
		\end{block}
		\begin{enumerate}[(A)]
			\item Performing vulnerability assessments
			\item Maintaining asset inventory
			\item Deploying changes without testing
			\item Conducting formal risk assessment
		\end{enumerate}
	\end{frame}
	
	\begin{frame}{Q18: Explanation}
		\begin{block}{Correct Answer: D}
			Conducting formal risk assessment is a requirement of many legal and regulatory compliance frameworks.
		\end{block}
		\begin{block}{Option Analysis}
			\begin{itemize}
				\item \textbf{A - Incorrect:} Vulnerability assessments are important but not universally mandatory
				\item \textbf{B - Incorrect:} Asset inventory is good practice but not always required
				\item \textbf{C - Incorrect:} This is poor practice
				\item \textbf{D - Correct:} Formal risk assessment is widely mandated
			\end{itemize}
		\end{block}
	\end{frame}
	
	\begin{frame}{Q19: Event vs Incident vs Breach}
		\begin{block}{Question}
			A healthcare employee accidentally sent more than 5,000 patient records in response to a phishing email. This is a:
		\end{block}
		\begin{enumerate}[(A)]
			\item Issue
			\item Event
			\item Incident
			\item Breach
		\end{enumerate}
	\end{frame}
	
	\begin{frame}{Q19: Explanation}
		\begin{block}{Correct Answer: D}
			This is a breach - unauthorized disclosure of PHI affecting more than 500 individuals.
		\end{block}
		\begin{block}{Option Analysis}
			\begin{itemize}
				\item \textbf{A - Incorrect:} Issue is a condition, not an occurrence
				\item \textbf{B - Incorrect:} Event alone doesn't cause negative outcome
				\item \textbf{C - Incorrect:} Incident is one or more events with negative outcome
				\item \textbf{D - Correct:} Breach is unauthorized disclosure
			\end{itemize}
		\end{block}
	\end{frame}
	
	\section{Domain 2: Threat, Vulnerability, Risk}
	

	
	\begin{frame}{Q22: CXXX}
		\begin{block}{Question}
			Which scoring technique is used to quantify a vulnerability?
		\end{block}
		\begin{enumerate}[(A)]
			\item CSS
			\item CSV
			\item CVSS
			\item CVVS
		\end{enumerate}
	\end{frame}
	
	\begin{frame}{Q22: Explanation}
		\begin{block}{Correct Answer: C}
			Common Vulnerability Scoring System (CVSS) is used to quantify vulnerabilities.
		\end{block}
		\begin{block}{Option Analysis}
			\begin{itemize}
				\item \textbf{A - Incorrect:} CSS is not vulnerability scoring
				\item \textbf{B - Incorrect:} CSV is file format
				\item \textbf{C - Correct:} CVSS is the standard
				\item \textbf{D - Incorrect:} CVVS is not standard
			\end{itemize}
		\end{block}
	\end{frame}
	
	
	
	\begin{frame}{Q22: Vulnerability Risk Assessment} \begin{block}{Question} An organization uses CVSS scores to prioritize vulnerability remediation. A vulnerability with a CVSS score of 9.8 affects a low-value internal system, while a vulnerability with a CVSS score of 6.5 affects a customer-facing payment application. What should the risk professional do FIRST? \end{block} \begin{enumerate}[(A)] \item Remediate the 9.8 vulnerability immediately because it has the highest score. \item Escalate both vulnerabilities directly to the board. \item Assess business impact and threat context before prioritizing remediation. \item Accept the 6.5 vulnerability because it is below 7.0. \end{enumerate} \end{frame} \begin{frame}{Q22: Explanation} \begin{block}{Correct Answer: C} CVSS measures technical severity, not business risk. Risk decisions should consider asset criticality, business impact, threat landscape, and organizational risk appetite in addition to CVSS scores. \end{block} \begin{block}{Option Analysis} \begin{itemize} \item \textbf{A - Incorrect:} Highest CVSS does not automatically mean highest business risk. \item \textbf{B - Incorrect:} Board escalation is not the first step. \item \textbf{C - Correct:} Business context must be evaluated before remediation priorities are established. \item \textbf{D - Incorrect:} Risk acceptance cannot be based solely on CVSS thresholds. \end{itemize} \end{block} \end{frame}
	
	
	\section{Domain 2: Risk Assessment Frameworks}
	
	\begin{frame}{Q23: Quantitative Risk Framework}
		\begin{block}{Question}
			Which framework is primarily used for quantitative risk management?
		\end{block}
		\begin{enumerate}[(A)]
			\item NIST 800-30
			\item FAIR
			\item ISO 27001
			\item ISO 27005
		\end{enumerate}
	\end{frame}
	
	\begin{frame}{Q23: Explanation}
		\begin{block}{Correct Answer: B}
			FAIR (Factor Analysis of Information Risk) is primarily used for quantitative risk management.
		\end{block}
		\begin{block}{Option Analysis}
			\begin{itemize}
				\item \textbf{A - Incorrect:} NIST 800-30 is qualitative
				\item \textbf{B - Correct:} FAIR is quantitative
				\item \textbf{C - Incorrect:} ISO 27001 is management standard
				\item \textbf{D - Incorrect:} ISO 27005 is qualitative
			\end{itemize}
		\end{block}
	\end{frame}
	
	
	\begin{frame}{Q23: Selecting a Risk Analysis Method} \begin{block}{Question} Senior management wants to compare the potential annual financial impact of ransomware risk against the cost of implementing additional security controls. Which approach would BEST support this analysis? \end{block} \begin{enumerate}[(A)] \item Use FAIR to estimate probable loss exposure in monetary terms. \item Use ISO 27001 to calculate annualized loss expectancy. \item Use NIST SP 800-30 to determine exact financial loss values. \item Use ISO 27005 to calculate return on security investment. \end{enumerate} \end{frame} \begin{frame}{Q23: Explanation} \begin{block}{Correct Answer: A} FAIR (Factor Analysis of Information Risk) is designed to quantify information risk in financial terms and supports cost-benefit and risk-based decision making. \end{block} \begin{block}{Option Analysis} \begin{itemize} \item \textbf{A - Correct:} FAIR provides quantitative estimates of probable loss exposure. \item \textbf{B - Incorrect:} ISO 27001 is an information security management standard, not a quantitative risk model. \item \textbf{C - Incorrect:} NIST SP 800-30 provides risk assessment guidance but does not produce precise financial loss calculations. \item \textbf{D - Incorrect:} ISO 27005 provides risk management guidance but is not primarily a quantitative financial analysis framework. \end{itemize} \end{block} \end{frame}
	

	
	\begin{frame}{Q26: Supply Chain Risk Framework}
		\begin{block}{Question}
			Which NIST framework provides guidance for supply chain management?
		\end{block}
		\begin{enumerate}[(A)]
			\item NIST 800-161
			\item NIST 800-30
			\item NIST 800-57
			\item ISO 27001
		\end{enumerate}
	\end{frame}
	
	\begin{frame}{Q26: Explanation}
		\begin{block}{Correct Answer: A}
			NIST SP 800-161 provides guidance for supply chain risk management.
		\end{block}
		\begin{block}{Option Analysis}
			\begin{itemize}
				\item \textbf{A - Correct:} 800-161 is supply chain
				\item \textbf{B - Incorrect:} 800-30 is risk management
				\item \textbf{C - Incorrect:} 800-57 is key management
				\item \textbf{D - Incorrect:} ISO 27001 is management standard
			\end{itemize}
		\end{block}
	\end{frame}
	
	\begin{frame}{Q27: Analysis}
		\begin{block}{Question}
			Which risk assessment technique provides results by displaying links between possible causes, controls, and consequences in a diagram?
		\end{block}
		\begin{enumerate}[(A)]
			\item FAIR
			\item BTA
			\item Markov analysis
			\item Monte Carlo analysis
		\end{enumerate}
	\end{frame}
	
	\begin{frame}{Q27: Explanation}
		\begin{block}{Correct Answer: B}
			Bow Tie Analysis (BTA) displays links between causes, controls, and consequences in a diagram.
		\end{block}
		\begin{block}{Option Analysis}
			\begin{itemize}
				\item \textbf{A - Incorrect:} FAIR is quantitative taxonomy
				\item \textbf{B - Correct:} BTA is diagram-based
				\item \textbf{C - Incorrect:} Markov uses state analysis
				\item \textbf{D - Incorrect:} Monte Carlo uses simulation
			\end{itemize}
		\end{block}
	\end{frame}
	
	
	\begin{frame}{Q27: Risk Analysis Technique} \begin{block}{Question} During a risk workshop, management wants to visualize how specific threat events could occur, the preventive controls in place, and the potential business consequences if those controls fail. Which risk assessment technique BEST meets this requirement? \end{block} \begin{enumerate}[(A)] \item FAIR \item Bow Tie Analysis (BTA) \item Markov Analysis \item Monte Carlo Simulation \end{enumerate} \end{frame} \begin{frame}{Q27: Explanation} \begin{block}{Correct Answer: B} Bow Tie Analysis (BTA) visually illustrates the relationship between threats, preventive controls, the risk event, recovery controls, and resulting consequences. \end{block} \begin{block}{Option Analysis} \begin{itemize} \item \textbf{A - Incorrect:} FAIR focuses on quantifying risk in financial terms. \item \textbf{B - Correct:} BTA uses a visual diagram linking causes, controls, events, and consequences. \item \textbf{C - Incorrect:} Markov analysis models state transitions and probabilities. \item \textbf{D - Incorrect:} Monte Carlo simulation evaluates a range of possible outcomes using repeated calculations. \end{itemize} \end{block} \end{frame}
	
	
	\begin{frame}{Q28: Risk Register Output}
		\begin{block}{Question}
			The results of a risk assessment should be summarized as a(n):
		\end{block}
		\begin{enumerate}[(A)]
			\item CAP
			\item Business continuity plan
			\item Organizational chart
			\item Risk register
		\end{enumerate}
	\end{frame}
	
	\begin{frame}{Q28: Explanation}
		\begin{block}{Correct Answer: D}
			Risk register is the immediate result of a risk assessment.
		\end{block}
		\begin{block}{Option Analysis}
			\begin{itemize}
				\item \textbf{A - Incorrect:} CAP is corrective action plan
				\item \textbf{B - Incorrect:} BCP is continuity planning
				\item \textbf{C - Incorrect:} Org chart shows structure
				\item \textbf{D - Correct:} Risk register documents risks
			\end{itemize}
		\end{block}
	\end{frame}
	
	\begin{frame}{Q28: Risk Assessment Deliverable} \begin{block}{Question} After completing an enterprise risk assessment, management wants a consolidated view of identified risks, their ratings, assigned owners, response plans, and monitoring status. Which document BEST fulfills this requirement? \end{block} \begin{enumerate}[(A)] \item Corrective Action Plan (CAP) \item Business Continuity Plan (BCP) \item Organizational Chart \item Risk Register \end{enumerate} \end{frame} \begin{frame}{Q28: Explanation} \begin{block}{Correct Answer: D} A risk register is the primary repository for documenting identified risks, risk ratings, ownership, response strategies, and ongoing monitoring activities. \end{block} \begin{block}{Option Analysis} \begin{itemize} \item \textbf{A - Incorrect:} A CAP addresses remediation actions for specific issues, not the overall inventory of risks. \item \textbf{B - Incorrect:} A BCP focuses on continuity and recovery from disruptions. \item \textbf{C - Incorrect:} An organizational chart defines reporting relationships. \item \textbf{D - Correct:} A risk register is the standard output of a risk assessment process. \end{itemize} \end{block} \end{frame}
	
	\section{Domain 2: Business Impact Analysis}
	
	\begin{frame}{Q29: BIA vs Risk Assessment}
		\begin{block}{Question}
			Which statement is false regarding risk assessment and BIA?
		\end{block}
		\begin{enumerate}[(A)]
			\item BIA identifies RPO, RTO, MTD
			\item Risk assessment identifies mitigation plans
			\item BCP needs BIA
			\item BIA and risk assessment are the same
		\end{enumerate}
	\end{frame}
	
	\begin{frame}{Q29: Explanation}
		\begin{block}{Correct Answer: D}
			BIA and risk assessment are NOT the same.
		\end{block}
		\begin{block}{Option Analysis}
			\begin{itemize}
				\item \textbf{A - True:} BIA identifies recovery objectives
				\item \textbf{B - True:} Risk assessment identifies mitigation
				\item \textbf{C - True:} BCP needs BIA as input
				\item \textbf{D - Not True:} They are different processes
			\end{itemize}
		\end{block}
	\end{frame}
	
	
	
	\begin{frame}{Q29: Business Impact Analysis vs. Risk Assessment} \begin{block}{Question} Which statement BEST describes the relationship between a Business Impact Analysis (BIA) and a risk assessment? \end{block} \begin{enumerate}[(A)] \item Both processes identify threats and vulnerabilities. \item A BIA determines business impact and recovery requirements, while a risk assessment evaluates risk likelihood and impact. \item A risk assessment establishes RTOs and RPOs, while a BIA identifies threat scenarios. \item Both processes produce identical outputs and can be used interchangeably. \end{enumerate} \end{frame} \begin{frame}{Q29: Explanation} \begin{block}{Correct Answer: B} A Business Impact Analysis determines the effect of disruptions on business operations and establishes recovery objectives. A risk assessment evaluates threats, vulnerabilities, likelihood, and impact to determine risk levels. \end{block} \begin{block}{Option Analysis} \begin{itemize} \item \textbf{A - Incorrect:} Threat and vulnerability identification is primarily part of risk assessment. \item \textbf{B - Correct:} BIA focuses on business impact and recovery requirements; risk assessment focuses on risk evaluation. \item \textbf{C - Incorrect:} RTO and RPO are outputs of a BIA. \item \textbf{D - Incorrect:} The two processes serve different purposes. \end{itemize} \end{block} \end{frame}
	
	
	\begin{frame}{Q30: Recovery Objectives}
		\begin{block}{Question}
			Which process is MOST likely to establish Recovery Time Objectives (RTOs)?
		\end{block}
		
		\begin{enumerate}[(A)]
			\item Risk Assessment
			\item Vulnerability Assessment
			\item Business Impact Analysis
			\item Control Self-Assessment
		\end{enumerate}
	\end{frame}
	
	\begin{frame}{Q30: Explanation}
		\begin{block}{Correct Answer: C}
			BIA establishes RTOs, RPOs, MTDs, and recovery priorities.
		\end{block}
		
		\begin{itemize}
			\item \textbf{A - Incorrect:} Risk assessment determines threats, vulnerabilities, likelihood, impact, and risk treatment. It does not establish recovery targets.
			\item \textbf{B - Incorrect:} Vulnerability assessments identify security weaknesses but do not determine business recovery objectives.
			\item \textbf{C - Correct:} BIA determines how quickly business services must be restored.
			\item \textbf{D - Incorrect:} Control Self-Assessment evaluates effectiveness of controls, not continuity requirements.
		\end{itemize}
	\end{frame}
	
	\begin{frame}{Q31: Threat Identification}
		\begin{block}{Question}
			Which activity is primarily focused on identifying threats and vulnerabilities?
		\end{block}
		
		\begin{enumerate}[(A)]
			\item Business Impact Analysis
			\item Risk Assessment
			\item Disaster Recovery Planning
			\item Business Continuity Testing
		\end{enumerate}
	\end{frame}
	
	\begin{frame}{Q31: Explanation}
		\begin{block}{Correct Answer: B}
			Risk assessment evaluates threats, vulnerabilities, likelihood, and impact.
		\end{block}
		
		\begin{itemize}
			\item \textbf{A - Incorrect:} BIA asks "what happens if a process fails?" not "what causes the failure?"
			\item \textbf{B - Correct:} Threat and vulnerability identification is a core risk assessment activity.
			\item \textbf{C - Incorrect:} DRP focuses on recovering technology after a disruption.
			\item \textbf{D - Incorrect:} Testing validates recovery plans already created.
		\end{itemize}
	\end{frame}
	
	\begin{frame}{Q32: Maximum Downtime}
		\begin{block}{Question}
			Maximum Tolerable Downtime (MTD) is determined during:
		\end{block}
		
		\begin{enumerate}[(A)]
			\item Risk Assessment
			\item Business Impact Analysis
			\item Penetration Testing
			\item Risk Monitoring
		\end{enumerate}
	\end{frame}
	
	\begin{frame}{Q32: Explanation}
		\begin{block}{Correct Answer: B}
			MTD is a business recovery metric determined during BIA.
		\end{block}
		
		\begin{itemize}
			\item \textbf{A - Incorrect:} Risk assessment determines risk levels rather than allowable downtime.
			\item \textbf{B - Correct:} BIA identifies how long a process can remain unavailable before unacceptable damage occurs.
			\item \textbf{C - Incorrect:} Penetration testing evaluates exploitability of vulnerabilities.
			\item \textbf{D - Incorrect:} Risk monitoring tracks changes in risks over time.
		\end{itemize}
	\end{frame}
	
	\begin{frame}{Q33: Residual Risk}
		\begin{block}{Question}
			Residual risk is MOST likely to be evaluated during:
		\end{block}
		
		\begin{enumerate}[(A)]
			\item Business Impact Analysis
			\item Business Continuity Planning
			\item Risk Assessment
			\item Disaster Recovery Testing
		\end{enumerate}
	\end{frame}
	
	\begin{frame}{Q33: Explanation}
		\begin{block}{Correct Answer: C}
			Residual risk is the risk remaining after controls operate.
		\end{block}
		
		\begin{itemize}
			\item \textbf{A - Incorrect:} BIA focuses on business impact and recovery requirements.
			\item \textbf{B - Incorrect:} BCP develops continuity strategies, not risk calculations.
			\item \textbf{C - Correct:} Residual risk is measured during risk assessment and risk evaluation.
			\item \textbf{D - Incorrect:} DR testing validates recovery capability.
		\end{itemize}
	\end{frame}
	
	\begin{frame}{Q34: Critical Processes}
		\begin{block}{Question}
			Which activity identifies business-critical processes and recovery priorities?
		\end{block}
		
		\begin{enumerate}[(A)]
			\item Risk Assessment
			\item Business Impact Analysis
			\item Control Assessment
			\item Audit Planning
		\end{enumerate}
	\end{frame}
	
	\begin{frame}{Q34: Explanation}
		\begin{block}{Correct Answer: B}
			BIA identifies critical processes and their recovery priorities.
		\end{block}
		
		\begin{itemize}
			\item \textbf{A - Incorrect:} Risk assessment identifies risks to processes, not which processes are most critical.
			\item \textbf{B - Correct:} BIA ranks processes based on business impact.
			\item \textbf{C - Incorrect:} Control assessment evaluates effectiveness of controls.
			\item \textbf{D - Incorrect:} Audit planning determines audit scope and resources.
		\end{itemize}
	\end{frame}
	
	\begin{frame}{Q35: Risk Response}
		\begin{block}{Question}
			Risk avoidance, mitigation, transfer, and acceptance are outputs of:
		\end{block}
		
		\begin{enumerate}[(A)]
			\item Risk Assessment
			\item Business Impact Analysis
			\item Business Continuity Planning
			\item Capacity Planning
		\end{enumerate}
	\end{frame}
	
	\begin{frame}{Q35: Explanation}
		\begin{block}{Correct Answer: A}
			These are the four classic risk response strategies.
		\end{block}
		
		\begin{itemize}
			\item \textbf{A - Correct:} Risk assessment leads to risk treatment decisions.
			\item \textbf{B - Incorrect:} BIA determines business recovery requirements.
			\item \textbf{C - Incorrect:} BCP focuses on operational continuity during disruptions.
			\item \textbf{D - Incorrect:} Capacity planning manages resource utilization and growth.
		\end{itemize}
	\end{frame}
	
	
	\begin{frame}{Q30: RPO Direction}
		\begin{block}{Question}
			Which is backward-looking in relation to BIA?
		\end{block}
		\begin{enumerate}[(A)]
			\item RPO
			\item RTO
			\item MTD
			\item SSO
		\end{enumerate}
	\end{frame}
	
	\begin{frame}{Q30: Explanation}
		\begin{block}{Correct Answer: A}
			RPO (Recovery Point Objective) is backward-looking - it measures acceptable data loss.
		\end{block}
		\begin{block}{Option Analysis}
			\begin{itemize}
				\item \textbf{A - Correct:} RPO is backward-looking
				\item \textbf{B - Incorrect:} RTO is forward-looking
				\item \textbf{C - Incorrect:} MTD is total downtime
				\item \textbf{D - Incorrect:} SSO is single sign-on
			\end{itemize}
		\end{block}
	\end{frame}
	
	
	\begin{frame}{Q30: RPO Direction} \begin{block}{Question} Which metric is considered backward-looking when determining recovery requirements? \end{block} \begin{enumerate}[(A)] \item Recovery Point Objective (RPO) \item Recovery Time Objective (RTO) \item Maximum Tolerable Downtime (MTD) \item Single Sign-On (SSO) \end{enumerate} \end{frame} \begin{frame}{Q30: Explanation} \begin{block}{Correct Answer: A} RPO looks backward from the disruption and defines the maximum acceptable amount of data loss measured in time. \end{block} \begin{itemize} \item \textbf{A - Correct:} RPO answers: ``How far back can we go and still recover without unacceptable data loss?'' \item \textbf{B - Incorrect:} RTO is forward-looking. It answers: ``How quickly must the system be restored after an outage?'' \item \textbf{C - Incorrect:} MTD represents the maximum time a business process can remain unavailable before severe business damage occurs. \item \textbf{D - Incorrect:} SSO (Single Sign-On) is an authentication technology and is unrelated to BIA recovery metrics. \end{itemize} \end{frame}
	

	
	\begin{frame}{Q32: XXX Risk}
		\begin{block}{Question}
			The risk after implementing controls in a real-time scenario against threats is called:
		\end{block}
		\begin{enumerate}[(A)]
			\item Residual risk
			\item Inherent risk
			\item Current risk
			\item Total risk
		\end{enumerate}
	\end{frame}
	
	\begin{frame}{Q32: Explanation}
		\begin{block}{Correct Answer: C}
			Current risk is the residual risk in a real-time threat scenario at a specific point in time.
		\end{block}
		\begin{block}{Option Analysis}
			\begin{itemize}
				\item \textbf{A - Incorrect:} Residual is not real-time
				\item \textbf{B - Incorrect:} Inherent is before controls
				\item \textbf{C - Correct:} Current is real-time residual risk
				\item \textbf{D - Incorrect:} Total risk is before controls
			\end{itemize}
		\end{block}
	\end{frame}
	
	% ================ SECTION 3: DOMAIN 3 - RISK RESPONSE ================
	
	\section{Domain 3: Risk Response}
	
	\begin{frame}{Q33: Risk Owner Purpose}
		\begin{block}{Question}
			The primary reason for having a risk owner is:
		\end{block}
		\begin{enumerate}[(A)]
			\item Leverage their resources
			\item Ensure accountability
			\item Provide flexibility
			\item Seek assistance
		\end{enumerate}
	\end{frame}
	
	\begin{frame}{Q33: Explanation}
		\begin{block}{Correct Answer: B}
			The risk owner drives accountability in risk remediation.
		\end{block}
		\begin{block}{Option Analysis}
			\begin{itemize}
				\item \textbf{A - Incorrect:} Resources are secondary
				\item \textbf{B - Correct:} Accountability is primary
				\item \textbf{C - Incorrect:} Flexibility is not primary
				\item \textbf{D - Incorrect:} Assistance is secondary
			\end{itemize}
		\end{block}
	\end{frame}
	
	
	\begin{frame}{Q33: Risk Owner Accountability} \begin{block}{Question} What is the PRIMARY reason for assigning a risk owner to an identified risk? \end{block} \begin{enumerate}[(A)] \item To leverage the risk owner's resources \item To ensure accountability for managing the risk \item To provide flexibility in risk response decisions \item To obtain assistance from subject matter experts \end{enumerate} \end{frame} \begin{frame}{Q33: Explanation} \begin{block}{Correct Answer: B} Risk ownership exists primarily to ensure that a specific individual is accountable for monitoring, responding to, and reporting on a risk. \end{block} \begin{itemize} \item \textbf{A - Incorrect:} A risk owner may have resources to address a risk, but resources are not why ownership is assigned. Ownership is about responsibility, not budget or staffing. \item \textbf{B - Correct:} Accountability is the foundation of effective risk governance. Without a clearly assigned owner, risks may not be monitored, mitigated, or escalated appropriately. \item \textbf{C - Incorrect:} Flexibility can exist regardless of ownership. Assigning an owner establishes responsibility rather than increasing flexibility. \item \textbf{D - Incorrect:} Subject matter experts may support the owner, but they do not replace accountability. A risk can have many contributors but should have one accountable owner. \end{itemize} \end{frame}
	
	\begin{frame}{Q34: NOT a Risk Response Strategy}
		\begin{block}{Question}
			Which is NOT a risk response strategy?
		\end{block}
		\begin{enumerate}[(A)]
			\item Accept
			\item Transfer
			\item Treat
			\item Mitigate
		\end{enumerate}
	\end{frame}
	
	\begin{frame}{Q34: Explanation}
		\begin{block}{Correct Answer: C}
			"Treat" is sometimes used interchangeably but is not a formal risk response strategy. The four strategies are Accept, Transfer, Mitigate, and Avoid.
		\end{block}
		\begin{block}{Option Analysis}
			\begin{itemize}
				\item \textbf{A - Incorrect:} Accept is a strategy
				\item \textbf{B - Incorrect:} Transfer is a strategy
				\item \textbf{C - Correct:} Treat is not a formal strategy
				\item \textbf{D - Incorrect:} Mitigate is a strategy
			\end{itemize}
		\end{block}
	\end{frame}
	
	\begin{frame}{Q35: Compensating Controls}
		\begin{block}{Question}
			A risk manager determines primary controls are insufficient and implements compensating controls. This is:
		\end{block}
		\begin{enumerate}[(A)]
			\item Risk mitigation
			\item Risk acceptance
			\item Risk sharing
			\item Risk avoidance
		\end{enumerate}
	\end{frame}
	
	\begin{frame}{Q35: Explanation}
		\begin{block}{Correct Answer: A}
			Implementing additional controls is an example of risk mitigation.
		\end{block}
		\begin{block}{Option Analysis}
			\begin{itemize}
				\item \textbf{A - Correct:} Adding controls is mitigation
				\item \textbf{B - Incorrect:} Acceptance would do nothing
				\item \textbf{C - Incorrect:} Sharing would involve third party
				\item \textbf{D - Incorrect:} Avoidance would terminate activity
			\end{itemize}
		\end{block}
	\end{frame}
	
	
	\begin{frame}{Q35: Compensating Controls} \begin{block}{Question} A risk manager determines existing controls are insufficient and implements compensating controls to reduce the likelihood of a risk event. Which risk response strategy is being applied? \end{block} \begin{enumerate}[(A)] \item Risk mitigation \item Risk acceptance \item Risk sharing \item Risk avoidance \end{enumerate} \end{frame} \begin{frame}{Q35: Explanation} \begin{block}{Correct Answer: A} Implementing compensating controls reduces either the likelihood or impact of a risk and is therefore a form of risk mitigation. \end{block} \begin{itemize} \item \textbf{A - Correct:} Risk mitigation reduces risk through controls, safeguards, monitoring, or process improvements. \item \textbf{B - Incorrect:} Risk acceptance means management knowingly accepts the residual risk without implementing additional controls. \item \textbf{C - Incorrect:} Risk sharing transfers part of the risk to another party, such as through insurance, outsourcing, or contractual agreements. \item \textbf{D - Incorrect:} Risk avoidance eliminates the activity giving rise to the risk altogether. \end{itemize} \end{frame}
	
	
	
	\begin{frame}{Q36: Insurance as Risk Response}
		\begin{block}{Question}
			Senior management buys insurance for an earthquake-prone site. This is:
		\end{block}
		\begin{enumerate}[(A)]
			\item Risk mitigation
			\item Risk acceptance
			\item Risk sharing
			\item Risk avoidance
		\end{enumerate}
	\end{frame}
	
	\begin{frame}{Q36: Explanation}
		\begin{block}{Correct Answer: C}
			Insurance transfers/sharing the risk with the insurance company.
		\end{block}
		\begin{block}{Option Analysis}
			\begin{itemize}
				\item \textbf{A - Incorrect:} Mitigation would reduce risk
				\item \textbf{B - Incorrect:} Acceptance would retain risk
				\item \textbf{C - Correct:} Sharing through insurance
				\item \textbf{D - Incorrect:} Avoidance would close site
			\end{itemize}
		\end{block}
	\end{frame}
	
	
	\begin{frame}{Q36: Insurance as a Risk Response} \begin{block}{Question} Senior management purchases insurance coverage for a data center located in an earthquake-prone region. Which risk response strategy is being applied? \end{block} \begin{enumerate}[(A)] \item Risk mitigation \item Risk acceptance \item Risk sharing \item Risk avoidance \end{enumerate} \end{frame} \begin{frame}{Q36: Explanation} \begin{block}{Correct Answer: C} Insurance transfers part of the financial consequences of a loss to a third party and is therefore considered risk sharing (or risk transfer). \end{block} \begin{itemize} \item \textbf{A - Incorrect:} Mitigation reduces the likelihood or impact of a risk through controls. Insurance does not prevent earthquakes or reduce their probability. \item \textbf{B - Incorrect:} Acceptance means management knowingly retains the risk without additional treatment. Purchasing insurance is a deliberate treatment strategy. \item \textbf{C - Correct:} Insurance shifts part of the financial burden to the insurer. The organization still experiences the event, but shares the loss. \item \textbf{D - Incorrect:} Avoidance eliminates the activity creating the risk. For example, relocating the data center away from the earthquake zone would be risk avoidance. \end{itemize} \end{frame}
	
	\begin{frame}{Q37: Risk What to do}
		\begin{block}{Question}
			Management terminates onboarding a new tool as it doesn't adapt to changing requirements. This is:
		\end{block}
		\begin{enumerate}[(A)]
			\item Risk mitigation
			\item Risk acceptance
			\item Risk sharing
			\item Risk avoidance
		\end{enumerate}
	\end{frame}
	
	\begin{frame}{Q37: Explanation}
		\begin{block}{Correct Answer: D}
			Terminating a project to eliminate risk is risk avoidance.
		\end{block}
		\begin{block}{Option Analysis}
			\begin{itemize}
				\item \textbf{A - Incorrect:} Mitigation would modify the tool
				\item \textbf{B - Incorrect:} Acceptance would proceed with risk
				\item \textbf{C - Incorrect:} Sharing would involve third party
				\item \textbf{D - Correct:} Avoidance eliminates the risk
			\end{itemize}
		\end{block}
	\end{frame}
	
	
	\begin{frame}{Q37: Product Modernization Decision} \begin{block}{Question} During implementation of a new enterprise platform, management determines the solution cannot support anticipated regulatory requirements and would require extensive future redevelopment. After reviewing the projected costs and risks, management decides not to proceed with implementation. Which risk response strategy has management selected? \end{block} \begin{enumerate}[(A)] \item Risk mitigation \item Risk acceptance \item Risk sharing \item Risk avoidance \end{enumerate} \end{frame} \begin{frame}{Q37: Explanation} \begin{block}{Correct Answer: D} By deciding not to proceed with the project, management eliminates the exposure associated with the initiative. \end{block} \begin{itemize} \item \textbf{A - Incorrect:} Mitigation would involve implementing controls or redesigning the solution while continuing the project. \item \textbf{B - Incorrect:} Acceptance would mean proceeding despite the identified concerns. \item \textbf{C - Incorrect:} Sharing would involve transferring part of the risk through contracts, insurance, or third parties. \item \textbf{D - Correct:} Cancelling the initiative eliminates the source of the risk and is therefore risk avoidance. \end{itemize} \end{frame}
	
	

	\begin{frame}{Q46: Risk Treatment Strategy} \begin{block}{Question} An organization identifies a concentration risk resulting from a critical business process being performed by a single employee. Management implements cross-training and succession procedures to reduce the likelihood of operational disruption. Which risk response strategy is MOST appropriate? \begin{enumerate}[(A)] \item Risk mitigation \item Risk acceptance \item Risk sharing \item Risk avoidance \end{enumerate} \textbf{Correct Answer: A} \vspace{0.2cm} \textbf{Explanation:} Cross-training and succession planning reduce the likelihood and impact of disruption while allowing the activity to continue. This is an example of risk mitigation. \vspace{0.2cm} \textbf{Option Analysis} \textbf{A - Correct: Risk Mitigation} Additional controls were implemented to reduce operational risk. The activity continues, but exposure is reduced. \vspace{0.2cm} \textbf{B - Incorrect: Risk Acceptance} Acceptance would mean management recognizes the risk but decides not to implement additional controls. \textit{Exam Clue: "Management formally accepts residual risk."} \vspace{0.2cm} \textbf{C - Incorrect: Risk Sharing} Sharing requires another party to assume some of the exposure through insurance, outsourcing, indemnification clauses, or contractual agreements. \textit{Exam Clue: "Transfer financial consequences to a third party."} \vspace{0.2cm} \textbf{D - Incorrect: Risk Avoidance} Avoidance would eliminate the activity entirely. For example, the organization may discontinue the process rather than improve it. \textit{Exam Clue: "Terminate, discontinue, withdraw, or cancel."} \vspace{0.2cm} \textbf{Why Candidates Miss This} Candidates often focus on the existence of risk and choose sharing or acceptance. The key phrase is that management implemented new controls while continuing the activity. In CRISC, that almost always indicates risk mitigation. \end{block} \end{frame}
	
	
	
	\section{Domain 3: Third-Party Risk Management}
	
	\begin{frame}{Q39: SLA Breach}
		\begin{block}{Question}
			Which would bind a third party to provide monetary credits in case of service failure?
		\end{block}
		\begin{enumerate}[(A)]
			\item Master service agreement
			\item Service-level agreement
			\item Non-disclosure agreement
			\item External audit
		\end{enumerate}
	\end{frame}
	
	\begin{frame}{Q39: Explanation}
		\begin{block}{Correct Answer: B}
			A breach in the SLA allows the organization to demand monetary credit.
		\end{block}
		\begin{block}{Option Analysis}
			\begin{itemize}
				\item \textbf{A - Incorrect:} MSA is general framework
				\item \textbf{B - Correct:} SLA defines service levels and penalties
				\item \textbf{C - Incorrect:} NDA protects confidentiality
				\item \textbf{D - Incorrect:} Audit is verification
			\end{itemize}
		\end{block}
	\end{frame}
	
	
\begin{frame}{Q39: Third-Party Service Risk}
	\begin{block}{Question}
		A critical application is hosted by a cloud provider. Management
		wants contractual recourse if system availability falls below
		agreed performance thresholds.
		
		Which document would MOST directly establish measurable uptime
		targets and associated service credits?
	\end{block}
	
	\begin{enumerate}[(A)]
		\item Master Service Agreement (MSA)
		\item Service-Level Agreement (SLA)
		\item Non-Disclosure Agreement (NDA)
		\item Independent Assurance Report
	\end{enumerate}
\end{frame}

\begin{frame}{Q39: Explanation}
	\begin{block}{Correct Answer: B - Service-Level Agreement (SLA)}
		An SLA defines measurable performance obligations such as uptime,
		response times, and service credits when agreed service levels are
		not achieved. When an organization seeks contractual remedies for
		poor service performance, the SLA is the document that provides
		the specific performance commitments and associated penalties.
	\end{block}
	
	\begin{block}{Option Analysis}
		
		\textbf{A - Incorrect: Master Service Agreement (MSA)}
		
		The MSA establishes the overall legal and commercial framework
		governing the relationship between the customer and the service
		provider. It commonly addresses liability, payment terms,
		governance requirements, dispute resolution procedures, and
		general contractual obligations. Candidates are often tempted
		to choose the MSA because it is the primary contract between the
		parties. However, detailed service-performance metrics, uptime
		guarantees, and service-credit provisions are normally addressed
		within the SLA rather than the MSA itself.
		
		\vspace{0.2cm}
		
		\textbf{B - Correct: Service-Level Agreement (SLA)}
		
		An SLA defines measurable operational expectations including
		availability targets, uptime requirements, response times,
		resolution targets, and service-credit provisions. When service
		levels are not achieved, the SLA specifies the remedies available
		to the customer. Because the question focuses on measurable
		performance commitments and monetary credits for service failures,
		the SLA is the best answer.
		
		\vspace{0.2cm}
		
		\textbf{C - Incorrect: Non-Disclosure Agreement (NDA)}
		
		An NDA focuses on confidentiality obligations and governs how
		sensitive information may be accessed, disclosed, used, stored,
		and protected. While it is an important third-party risk
		management document, it does not establish availability targets,
		uptime commitments, performance guarantees, or service-credit
		provisions.
		
		\vspace{0.2cm}
		
		\textbf{D - Incorrect: Independent Assurance Report}
		
		Independent assurance reports such as SOC reports or external
		audit reports provide evidence regarding the design and operating
		effectiveness of controls. These reports help management assess
		third-party risk, but they do not create legally enforceable
		performance obligations. Assurance reports measure and validate
		controls, whereas SLAs define contractual performance
		expectations and penalties.
		
	\end{block}
	
	\begin{block}{CRISC Exam Tip}
		Remember the distinction: the MSA governs the overall business
		relationship, the SLA governs measurable performance obligations,
		the NDA governs confidentiality requirements, and assurance
		reports provide evidence that controls are functioning as
		expected. If the question mentions uptime, availability,
		performance metrics, penalties, or service credits, the answer is
		usually the SLA.
	\end{block}
\end{frame}
	
	\begin{frame}{Q40: NDA Purpose}
		\begin{block}{Question}
			Which should be signed with third party to protect intellectual property?
		\end{block}
		\begin{enumerate}[(A)]
			\item Master service agreement
			\item Service-level agreement
			\item Non-disclosure agreement
			\item External audit
		\end{enumerate}
	\end{frame}
	
	\begin{frame}{Q40: Explanation}
		\begin{block}{Correct Answer: C}
			NDAs protect intellectual property and interests of the organization.
		\end{block}
		\begin{block}{Option Analysis}
			\begin{itemize}
				\item \textbf{A - Incorrect:} MSA is general framework
				\item \textbf{B - Incorrect:} SLA defines service levels
				\item \textbf{C - Correct:} NDA protects IP and confidentiality
				\item \textbf{D - Incorrect:} Audit is verification
			\end{itemize}
		\end{block}
	\end{frame}
	
	\begin{frame}{Q41: External Audit Output}
		\begin{block}{Question}
			Which is NOT the final output of an external audit?
		\end{block}
		\begin{enumerate}[(A)]
			\item SOC 2 report
			\item ISO 27001 certification
			\item HITRUST certification
			\item Non-disclosure agreement
		\end{enumerate}
	\end{frame}
	
	\begin{frame}{Q41: Explanation}
		\begin{block}{Correct Answer: D}
			Non-disclosure agreements are signed internally, not the output of an external audit.
		\end{block}
		\begin{block}{Option Analysis}
			\begin{itemize}
				\item \textbf{A - Incorrect:} SOC 2 is audit output
				\item \textbf{B - Incorrect:} ISO 27001 is audit output
				\item \textbf{C - Incorrect:} HITRUST is audit output
				\item \textbf{D - Correct:} NDA is not an audit output
			\end{itemize}
		\end{block}
	\end{frame}
	
	
	\begin{frame}{Q41: Third-Party Assurance} \begin{block}{Question} A bank is assessing a critical cloud provider and wants independent evidence that security controls operate effectively. Which would LEAST help satisfy this requirement? \end{block} \begin{enumerate}[(A)] \item SOC 2 Type II report \item ISO 27001 certification \item HITRUST certification \item Non-Disclosure Agreement (NDA) \end{enumerate} \end{frame} 
	
	
	
	\begin{frame}[shrink=10]{Q41: Explanation} \begin{block}{Correct Answer: D} An NDA protects confidential information but does not provide assurance that security controls are designed or operating effectively. \end{block} \begin{block}{Option Analysis} \textbf{A - Incorrect:} A SOC 2 Type II report provides independent assurance regarding the effectiveness of controls. \vspace{0.15cm} \textbf{B - Incorrect:} ISO 27001 certification demonstrates that an information security management system has been assessed against a recognized standard. \vspace{0.15cm} \textbf{C - Incorrect:} HITRUST certification provides independent validation of security and compliance controls. \vspace{0.15cm} \textbf{D - Correct:} An NDA governs confidentiality obligations but does not verify the existence or effectiveness of controls. \end{block} \begin{block}{CRISC Tip} Assurance = SOC, ISO, HITRUST. Confidentiality = NDA. \end{block} \end{frame}
	
	\begin{frame}{Q42: Exception Review Frequency}
		\begin{block}{Question}
			The risk practitioner should review and verify granted exceptions are still required at least:
		\end{block}
		\begin{enumerate}[(A)]
			\item Weekly
			\item Monthly
			\item Quarterly
			\item Annually
		\end{enumerate}
	\end{frame}
	
	\begin{frame}{Q42: Explanation}
		\begin{block}{Correct Answer: D}
			Granted exceptions should be verified at least annually.
		\end{block}
		\begin{block}{Option Analysis}
			\begin{itemize}
				\item \textbf{A - Incorrect:} Weekly is too frequent
				\item \textbf{B - Incorrect:} Monthly is too frequent
				\item \textbf{C - Incorrect:} Quarterly is too frequent
				\item \textbf{D - Correct:} Annual review is standard
			\end{itemize}
		\end{block}
		
		
		\begin{block}{CRISC Exam Tip} Whenever you see "exception," think: approval, documentation, compensating controls, risk owner, expiration date, and periodic review. \end{block}
		
	\end{frame}
	
	
	\begin{frame}{Q42: Policy Exception Governance} \begin{block}{Question} A business unit was granted an exception to the organization's encryption standard because a legacy application could not support the required controls. The compensating controls remain in place, and the exception was approved by management. What should the risk practitioner do NEXT? \end{block} \begin{enumerate}[(A)] \item Close the exception because management approved it \item Review the exception periodically to ensure it remains justified \item Transfer ownership of the exception to the vendor \item Remove the compensating controls \end{enumerate} \end{frame} 
	
	
	\begin{frame}[shrink=10]{Q42: Explanation} \begin{block}{Correct Answer: B} Approved exceptions should be periodically reviewed to confirm they remain necessary, justified, and within the organization's risk appetite. \end{block} \begin{block}{Option Analysis} \textbf{A - Incorrect:} Management approval does not eliminate the need for ongoing oversight. Risk exceptions require periodic review because business conditions, technology, and risk exposure can change. \vspace{0.15cm} \textbf{B - Correct:} Exception governance requires regular reassessment to determine whether the original justification remains valid and whether the exception should be renewed, modified, or terminated. \vspace{0.15cm} \textbf{C - Incorrect:} Risk ownership and accountability remain with the organization even when third parties are involved. \vspace{0.15cm} \textbf{D - Incorrect:} Compensating controls were implemented to reduce risk exposure. Removing them would increase residual risk. \end{block} 
	 \end{frame}
	
	\begin{frame}{Q43: CAB Composition}
		\begin{block}{Question}
			The CAB should consist of which of the following?
		\end{block}
		\begin{enumerate}[(A)]
			\item Only IT
			\item IT and engineering
			\item IT, engineering, and security
			\item All relevant stakeholders
		\end{enumerate}
	\end{frame}
	
	\begin{frame}{Q43: Explanation}
		\begin{block}{Correct Answer: D}
			CAB should consist of all relevant stakeholders.
		\end{block}
		\begin{block}{Option Analysis}
			\begin{itemize}
				\item \textbf{A - Incorrect:} IT alone is insufficient
				\item \textbf{B - Incorrect:} Engineering alone is insufficient
				\item \textbf{C - Incorrect:} Security is needed but not sufficient
				\item \textbf{D - Correct:} All relevant stakeholders needed
			\end{itemize}
		\end{block}
	\end{frame}
	
	\section{Domain 3: Control Design}
	
	\begin{frame}{Q44: Control Categories}
		\begin{block}{Question}
			Which is NOT a control category?
		\end{block}
		\begin{enumerate}[(A)]
			\item Additive
			\item Preventive
			\item Deterrent
			\item Detective
		\end{enumerate}
	\end{frame}
	
	\begin{frame}{Q44: Explanation}
		\begin{block}{Correct Answer: A}
			Additive is not a control category. The categories are Preventive, Detective, Corrective, Deterrent, Compensating.
		\end{block}
		\begin{block}{Option Analysis}
			\begin{itemize}
				\item \textbf{A - Correct:} Additive is not a control category
				\item \textbf{B - Incorrect:} Preventive is a category
				\item \textbf{C - Incorrect:} Deterrent is a category
				\item \textbf{D - Incorrect:} Detective is a category
			\end{itemize}
		\end{block}
	\end{frame}
	
	
	\begin{frame}{Q44: Control Classification} \begin{block}{Question} A financial institution implements multi-factor authentication to reduce the likelihood of unauthorized access to customer accounts. How should this control BEST be classified? \end{block} \begin{enumerate}[(A)] \item Detective \item Corrective \item Preventive \item Compensating \end{enumerate} \end{frame} 
	
	\begin{frame}[shrink=10]{Q44: Explanation} \begin{block}{Correct Answer: C} Multi-factor authentication attempts to stop unauthorized access before it occurs. Controls that reduce the likelihood of an event are preventive controls. \end{block} \begin{block}{Option Analysis} \text{A - Incorrect: Detective} Detective controls identify events that have already occurred. Examples include log reviews, security monitoring, intrusion detection systems, and audit trails. These controls provide visibility but do not stop the initial event. \textbf{B - Incorrect: Corrective} Corrective controls restore systems after an incident has occurred. Examples include backup restoration, incident response activities, and disaster recovery procedures.  \text{C - Correct: Preventive} Preventive controls attempt to stop undesirable events before they happen. Examples include MFA, segregation of duties, encryption, firewalls, and access controls. \vspace{0.15cm} \text{D - Incorrect: Compensating} Compensating controls are alternative controls used when a primary control cannot be implemented. For example, increased monitoring may compensate for a system that cannot support MFA. \end{block}  \end{frame}
	
	\begin{frame}{Q45: Proactive Control}
		\begin{block}{Question}
			Implementing controls proactively based on root cause analysis is an example of:
		\end{block}
		\begin{enumerate}[(A)]
			\item Preventive
			\item Detective
			\item Corrective
			\item Compensating
		\end{enumerate}
	\end{frame}
	
	\begin{frame}{Q45: Explanation}
		\begin{block}{Correct Answer: A}
			Implementing controls proactively is an example of preventive control.
		\end{block}
		\begin{block}{Option Analysis}
			\begin{itemize}
				\item \textbf{A - Correct:} Preventive stops issues before they occur
				\item \textbf{B - Incorrect:} Detective finds issues after
				\item \textbf{C - Incorrect:} Corrective fixes issues after
				\item \textbf{D - Incorrect:} Compensating compensates for weaknesses
			\end{itemize}
		\end{block}
	\end{frame}
	
	\begin{frame}{Q46: Which Control}
		\begin{block}{Question}
			Reviewing audit logs in a SIEM tool is an example of:
		\end{block}
		\begin{enumerate}[(A)]
			\item Corrective control
			\item Detective control
			\item Compensating control
			\item Deterrent control
		\end{enumerate}
	\end{frame}
	
	\begin{frame}{Q46: Explanation}
		\begin{block}{Correct Answer: B}
			Reviewing audit logs is detective control - it identifies issues after they occur.
		\end{block}
		\begin{block}{Option Analysis}
			\begin{itemize}
				\item \textbf{A - Incorrect:} Corrective fixes issues
				\item \textbf{B - Correct:} Detective identifies issues
				\item \textbf{C - Incorrect:} Compensating compensates
				\item \textbf{D - Incorrect:} Deterrent discourages actions
			\end{itemize}
		\end{block}
	\end{frame}
	
	
	
	
	\begin{frame}{Q46: Security Monitoring} \begin{block}{Question} During a security review, an analyst uses a SIEM platform to identify unusual login activity and investigate potential unauthorized access attempts. Which type of control is MOST evident? \end{block} \begin{enumerate}[(A)] \item Corrective control \item Detective control \item Compensating control \item Deterrent control \end{enumerate} 
	\begin{block}{CRISC Exam Tip} Preventive = Stop Detective = Discover Corrective = Recover Deterrent = Discourage Compensating = Alternative \end{block}
\end{frame}

	 \begin{frame}[shrink=10]{Q46: Explanation} \begin{block}{Correct Answer: B - Detective Control} Reviewing security events and audit logs allows an organization to identify suspicious activity after it occurs. The primary purpose is detection rather than prevention or recovery. \end{block} \begin{block}{Option Analysis} \textbf{A - Incorrect: Corrective Control} Corrective controls are used after an incident has been detected to restore normal operations or reduce damage. Examples include restoring backups, rebuilding systems, and incident-response activities. A SIEM helps discover problems but does not correct them. \vspace{0.15cm} \textbf{B - Correct: Detective Control} SIEM solutions collect and analyze logs to identify suspicious behavior, policy violations, and security incidents. The objective is to detect events that have already occurred or are currently occurring. \vspace{0.15cm} \textbf{C - Incorrect: Compensating Control} Compensating controls serve as alternatives when a primary control cannot be implemented. For example, increased monitoring may compensate for the inability to deploy a stronger preventive control. In this scenario, the question focuses on log review itself, which is detective in nature. \vspace{0.15cm} \textbf{D - Incorrect: Deterrent Control} Deterrent controls discourage undesirable behavior. Examples include warning banners, sanctions policies, surveillance notices, and visible security presence. These controls influence behavior but do not identify actual incidents. \end{block}  \end{frame}
	
	
	\begin{frame}{Q47: Deterrent Control}
		\begin{block}{Question}
			Installing security cameras to secure a data center is:
		\end{block}
		\begin{enumerate}[(A)]
			\item Corrective control
			\item Detective control
			\item Compensating control
			\item Deterrent control
		\end{enumerate}
	\end{frame}
	
	\begin{frame}{Q47: Explanation}
		\begin{block}{Correct Answer: D}
			Security cameras are a deterrent control - they discourage malicious activity.
		\end{block}
		\begin{block}{Option Analysis}
			\begin{itemize}
				\item \textbf{A - Incorrect:} Corrective would fix issues
				\item \textbf{B - Incorrect:} Detective would identify issues
				\item \textbf{C - Incorrect:} Compensating would compensate
				\item \textbf{D - Correct:} Deterrent discourages activity
			\end{itemize}
		\end{block}
	\end{frame}
	
	\begin{frame}{Q48: Which Control}
		\begin{block}{Question}
			A risk manager introduced additional oversight of an accounting system where logical controls are not implemented. This is:
		\end{block}
		\begin{enumerate}[(A)]
			\item Corrective control
			\item Detective control
			\item Compensating control
			\item Deterrent control
		\end{enumerate}
	\end{frame}
	
	\begin{frame}{Q48: Explanation}
		\begin{block}{Correct Answer: C}
			Additional oversight compensating for missing logical controls is a compensating control.
		\end{block}
		\begin{block}{Option Analysis}
			\begin{itemize}
				\item \textbf{A - Incorrect:} Corrective would fix the issue
				\item \textbf{B - Incorrect:} Detective would identify
				\item \textbf{C - Correct:} Compensating compensates for weakness
				\item \textbf{D - Incorrect:} Deterrent would discourage
			\end{itemize}
		\end{block}
	\end{frame}
	
	\begin{frame}{Q49: Which Control}
		\begin{block}{Question}
			An IT system failed during changes and lost data. Risk manager advises restoring from backups. This is:
		\end{block}
		\begin{enumerate}[(A)]
			\item Corrective control
			\item Detective control
			\item Compensating control
			\item Deterrent control
		\end{enumerate}
	\end{frame}
	
	\begin{frame}{Q49: Explanation}
		\begin{block}{Correct Answer: A}
			Restoring from backups to fix an issue is a corrective control.
		\end{block}
		\begin{block}{Option Analysis}
			\begin{itemize}
				\item \textbf{A - Correct:} Corrective fixes the issue
				\item \textbf{B - Incorrect:} Detective would identify
				\item \textbf{C - Incorrect:} Compensating would compensate
				\item \textbf{D - Incorrect:} Deterrent would discourage
			\end{itemize}
		\end{block}
	\end{frame}
	
	\begin{frame}{Q50: Post-Incident Controls}
		\begin{block}{Question}
			Which control categories are used after the incident has happened?
		\end{block}
		\begin{enumerate}[(A)]
			\item Corrective, compensating
			\item Corrective, deterrent
			\item Corrective, detective
			\item Deterrent, preventive
		\end{enumerate}
	\end{frame}
	
	\begin{frame}{Q50: Explanation}
		\begin{block}{Correct Answer: C}
			Detective and corrective controls are used after the incident has happened.
		\end{block}
		\begin{block}{Option Analysis}
			\begin{itemize}
				\item \textbf{A - Incorrect:} Compensating is ongoing
				\item \textbf{B - Incorrect:} Deterrent is before
				\item \textbf{C - Correct:} Detective and corrective are after
				\item \textbf{D - Incorrect:} Deterrent/preventive are before
			\end{itemize}
		\end{block}
	\end{frame}
	
	\begin{frame}{Q51: Parallel Changeover}
		\begin{block}{Question}
			Risk manager overseeing critical system upgrade with no downtime. Recommendation should be:
		\end{block}
		\begin{enumerate}[(A)]
			\item Parallel
			\item Phased
			\item Abrupt
			\item Any of the above
		\end{enumerate}
	\end{frame}
	
	\begin{frame}{Q51: Explanation}
		\begin{block}{Correct Answer: A}
			Parallel changeover allows both systems to run simultaneously with no downtime.
		\end{block}
		\begin{block}{Option Analysis}
			\begin{itemize}
				\item \textbf{A - Correct:} Parallel ensures no downtime
				\item \textbf{B - Incorrect:} Phased may have downtime
				\item \textbf{C - Incorrect:} Abrupt has high risk
				\item \textbf{D - Incorrect:} Parallel is best for no downtime
			\end{itemize}
		\end{block}
	\end{frame}
	
	\begin{frame}{Q52:  Changeover}
		\begin{block}{Question}
			System upgrading per module with no dependencies. Most efficient method?
		\end{block}
		\begin{enumerate}[(A)]
			\item Parallel
			\item Phased
			\item Abrupt
			\item Any of the above
		\end{enumerate}
	\end{frame}
	
	\begin{frame}{Q52: Explanation}
		\begin{block}{Correct Answer: B}
			Phased changeover is most efficient when modules are independent.
		\end{block}
		\begin{block}{Option Analysis}
			\begin{itemize}
				\item \textbf{A - Incorrect:} Parallel is expensive
				\item \textbf{B - Correct:} Phased works well with independent modules
				\item \textbf{C - Incorrect:} Abrupt is too risky
				\item \textbf{D - Incorrect:} Phased is most efficient
			\end{itemize}
		\end{block}
	\end{frame}
	
	
	\begin{frame}{Q52: System Implementation Strategy} \begin{block}{Question} A bank is replacing a legacy application with a modular solution. Each module can operate independently and can be implemented without affecting the others. Which implementation approach would MOST efficiently reduce deployment risk while allowing incremental rollout? \end{block} \begin{enumerate}[(A)] \item Parallel changeover \item Phased changeover \item Abrupt changeover \item Pilot changeover \end{enumerate} \end{frame} \begin{frame}[shrink=10]{Q52: Explanation} \footnotesize \begin{block}{Correct Answer: B - Phased Changeover} Phased implementation allows independent modules to be deployed incrementally. This reduces implementation risk while avoiding the cost and complexity of operating both systems simultaneously. \end{block} \begin{block}{Option Analysis} \textbf{A - Incorrect: Parallel Changeover} Parallel implementation operates the old and new systems simultaneously. While it provides strong fallback capability, it is resource-intensive and generally more expensive than necessary when modules can be migrated independently. \vspace{0.15cm} \textbf{B - Correct: Phased Changeover} Phased changeover introduces functionality gradually. Independent modules can be implemented sequentially, allowing issues to be isolated and resolved with reduced operational impact. \vspace{0.15cm} \textbf{C - Incorrect: Abrupt Changeover} Abrupt (or direct) changeover replaces the entire system at once. Although it may be faster, it introduces greater implementation risk because rollback options are limited. \vspace{0.15cm} \textbf{D - Incorrect: Pilot Changeover} A pilot approach deploys the new system to a limited group of users or locations before broader implementation. This reduces risk but does not specifically leverage the independence of individual modules. \end{block} \end{frame}
	
	\begin{frame}{Q53: Post-Implementation Audience}
		\begin{block}{Question}
			Who should be the target audience for a post-implementation review?
		\end{block}
		\begin{enumerate}[(A)]
			\item Developers
			\item Risk manager
			\item Risk owner
			\item All major stakeholders
		\end{enumerate}
	\end{frame}
	
	\begin{frame}{Q53: Explanation}
		\begin{block}{Correct Answer: D}
			All major stakeholders should participate in post-implementation review.
		\end{block}
		\begin{block}{Option Analysis}
			\begin{itemize}
				\item \textbf{A - Incorrect:} Developers alone insufficient
				\item \textbf{B - Incorrect:} Risk manager alone insufficient
				\item \textbf{C - Incorrect:} Risk owner alone insufficient
				\item \textbf{D - Correct:} All stakeholders needed
			\end{itemize}
		\end{block}
	\end{frame}
	
	\section{Domain 3: Monitoring and Reporting}
	
	\begin{frame}{Q54: Log Aggregation Purpose}
		\begin{block}{Question}
			What is the primary purpose of log aggregation in an organization?
		\end{block}
		\begin{enumerate}[(A)]
			\item Identify and analyze important events
			\item Synchronize clocks
			\item Eliminate need for logging
			\item Provide centralized storage only
		\end{enumerate}
	\end{frame}
	
	\begin{frame}{Q54: Explanation}
		\begin{block}{Correct Answer: A}
			Log aggregation helps identify and analyze important events and potential incidents.
		\end{block}
		\begin{block}{Option Analysis}
			\begin{itemize}
				\item \textbf{A - Correct:} Identification and analysis of events
				\item \textbf{B - Incorrect:} Clock sync is separate
				\item \textbf{C - Incorrect:} Doesn't eliminate logging
				\item \textbf{D - Incorrect:} Storage is not primary purpose
			\end{itemize}
		\end{block}
	\end{frame}
	
	\begin{frame}{Q55: Log Tool}
		\begin{block}{Question}
			Which is a commonly used tool for log aggregation and analysis?
		\end{block}
		\begin{enumerate}[(A)]
			\item SIEM system
			\item Vulnerability assessment tool
			\item Penetration testing tool
			\item Heat map generator
		\end{enumerate}
	\end{frame}
	
	\begin{frame}{Q55: Explanation}
		\begin{block}{Correct Answer: A}
			SIEM systems are commonly used for log aggregation and analysis.
		\end{block}
		\begin{block}{Option Analysis}
			\begin{itemize}
				\item \textbf{A - Correct:} SIEM handles log aggregation
				\item \textbf{B - Incorrect:} VA tools find vulnerabilities
				\item \textbf{C - Incorrect:} Pen testing finds exploits
				\item \textbf{D - Incorrect:} Heat maps are visualization
			\end{itemize}
		\end{block}
	\end{frame}
	
	
	\begin{frame}{Q55: Security Monitoring} \begin{block}{Question} During a security incident investigation, analysts need to collect logs from firewalls, servers, databases, and applications into a centralized platform where events can be correlated and analyzed. Which technology BEST supports this requirement? \end{block} \begin{enumerate}[(A)] \item Security Information and Event Management (SIEM) \item Vulnerability Assessment Tool \item Penetration Testing Tool \item Risk Heat Map \end{enumerate}
	
	
	\begin{block}{CRISC Exam Tip} Think about the primary purpose: SIEM = Log Collection and Monitoring Vulnerability Scanner = Find Weaknesses Penetration Testing = Exploit Weaknesses Heat Map = Visualize Risk \end{block}
	
	 \end{frame} \begin{frame}[shrink=10]{Q55: Explanation} \footnotesize \begin{block}{Correct Answer: A - Security Information and Event Management (SIEM)} SIEM platforms aggregate, normalize, correlate, and analyze logs from multiple sources, providing organizations with centralized monitoring and incident-detection capabilities. \end{block} \begin{block}{Option Analysis} \textbf{A - Correct: Security Information and Event Management (SIEM)} SIEM solutions collect logs from numerous systems and correlate events to identify suspicious activity, security incidents, and policy violations. This makes SIEM the preferred solution for centralized log monitoring and analysis. \vspace{0.15cm} \textbf{B - Incorrect: Vulnerability Assessment Tool} Vulnerability assessment tools identify known security weaknesses and missing patches. Although they help identify security risks, they do not primarily function as centralized log management platforms. \vspace{0.15cm} \textbf{C - Incorrect: Penetration Testing Tool} Penetration-testing tools are used to simulate attacks and identify exploitable weaknesses. Their purpose is to evaluate security posture rather than continuously aggregate and analyze operational logs. \vspace{0.15cm} \textbf{D - Incorrect: Risk Heat Map} A risk heat map is a reporting and visualization technique used to communicate risk levels based on likelihood and impact. It may use information derived from security events but does not perform log collection or event correlation. \end{block}  \end{frame}
	
	
	
	\begin{frame}{Q56: KRI Role}
		\begin{block}{Question}
			What is the role of KRIs in risk monitoring?
		\end{block}
		\begin{enumerate}[(A)]
			\item Predict risks that may breach thresholds
			\item Measure effectiveness of controls
			\item Indicate weaknesses in controls
			\item Assess control performance
		\end{enumerate}
	\end{frame}
	
	\begin{frame}{Q56: Explanation}
		\begin{block}{Correct Answer: A}
			KRIs predict risks that may breach predefined thresholds.
		\end{block}
		\begin{block}{Option Analysis}
			\begin{itemize}
				\item \textbf{A - Correct:} KRIs predict risk breaches
				\item \textbf{B - Incorrect:} KPIs measure effectiveness
				\item \textbf{C - Incorrect:} KCIs indicate weaknesses
				\item \textbf{D - Incorrect:} KPIs assess performance
			\end{itemize}
		\end{block}
	\end{frame}
	
	
	
	
	\begin{frame}{Q56: Risk Monitoring Metrics} \begin{block}{Question} A bank's risk committee receives a report showing that the number of failed privileged-access login attempts has increased by 40\% over the past quarter. Management is concerned that the trend may indicate increasing cybersecurity risk and could eventually exceed approved risk tolerance levels. Which metric is MOST likely being used? \end{block} \begin{enumerate}[(A)] \item Key Risk Indicator (KRI) \item Key Performance Indicator (KPI) \item Key Control Indicator (KCI) \item Key Audit Indicator (KAI) \end{enumerate}
	\begin{block}{CRISC Exam Tip} KRI = Is risk increasing? KPI = Is the business performing? KCI = Are controls working? \end{block}
	 \end{frame}
	
	
	 \begin{frame}[shrink=10]{Q56: Explanation} \footnotesize \begin{block}{Correct Answer: A - Key Risk Indicator (KRI)} KRIs provide early warning signals that risk exposure may be increasing and could eventually exceed established thresholds or risk appetite. \end{block} \begin{block}{Option Analysis} \textbf{A - Correct: Key Risk Indicator (KRI)} KRIs are forward-looking measures designed to identify changes in risk exposure. An increasing number of failed privileged-access attempts may indicate growing cyber threats and warrant additional monitoring or management attention. \vspace{0.15cm} \textbf{B - Incorrect: Key Performance Indicator (KPI)} KPIs measure business or operational performance. Examples include transaction-processing time, customer satisfaction, or system availability. The scenario focuses on increasing risk exposure rather than operational performance. \vspace{0.15cm} \textbf{C - Incorrect: Key Control Indicator (KCI)} KCIs measure whether controls are operating effectively. Examples include patch compliance rates, MFA enrollment percentages, or completion of access reviews. The scenario measures risk trends, not control effectiveness. \vspace{0.15cm} \textbf{D - Incorrect: Key Audit Indicator (KAI)} This is a distractor. CRISC commonly discusses KRIs, KPIs, and KCIs. The metric described in the scenario is specifically related to monitoring risk exposure. \end{block}  \end{frame}
	
	\begin{frame}{Q57: Dashboard Characteristics}
		\begin{block}{Question}
			Which format of risk reporting is known for flexibility in combining qualitative and quantitative metrics?
		\end{block}
		\begin{enumerate}[(A)]
			\item Executive summary
			\item Heat map
			\item Scorecard
			\item Dashboard
		\end{enumerate}
	\end{frame}
	
	\begin{frame}{Q57: Explanation}
		\begin{block}{Correct Answer: D}
			Dashboards combine qualitative and quantitative metrics with flexibility.
		\end{block}
		\begin{block}{Option Analysis}
			\begin{itemize}
				\item \textbf{A - Incorrect:} Summary is text-based
				\item \textbf{B - Incorrect:} Heat map is qualitative only
				\item \textbf{C - Incorrect:} Scorecard is qualitative
				\item \textbf{D - Correct:} Dashboard is flexible
			\end{itemize}
		\end{block}
	\end{frame}
	
	\begin{frame}{Q58: SMART Metrics}
		\begin{block}{Question}
			What are the key characteristics of a SMART metric?
		\end{block}
		\begin{enumerate}[(A)]
			\item Specific, Measurable, Attainable, Relevant, Timely
			\item Standardized, Meaningful, Achievable, Realistic, Timeless
			\item Secure, Measurable, Aligned, Reliable, Trustworthy
			\item Significant, Measurable, Accurate, Relevant, Timed
		\end{enumerate}
	\end{frame}
	
	\begin{frame}{Q58: Explanation}
		\begin{block}{Correct Answer: A}
			SMART stands for Specific, Measurable, Attainable, Relevant, Timely.
		\end{block}
		\begin{block}{Option Analysis}
			\begin{itemize}
				\item \textbf{A - Correct:} Standard SMART definition
				\item \textbf{B - Incorrect:} Not standard SMART
				\item \textbf{C - Incorrect:} Not standard SMART
				\item \textbf{D - Incorrect:} Not standard SMART
			\end{itemize}
		\end{block}
	\end{frame}
	
	\begin{frame}{Q59: Attack Testing}
		\begin{block}{Question}
			Which technique involves simulating a real attack to identify undiscovered vulnerabilities?
		\end{block}
		\begin{enumerate}[(A)]
			\item Self-assessment
			\item Internal audit
			\item Vulnerability assessment
			\item Penetration testing
		\end{enumerate}
	\end{frame}
	
	\begin{frame}{Q59: Explanation}
		\begin{block}{Correct Answer: D}
			Penetration testing simulates real attacks to identify vulnerabilities.
		\end{block}
		\begin{block}{Option Analysis}
			\begin{itemize}
				\item \textbf{A - Incorrect:} Self-assessment is internal review
				\item \textbf{B - Incorrect:} Audit is compliance review
				\item \textbf{C - Incorrect:} VA finds known vulnerabilities
				\item \textbf{D - Correct:} Penetration test simulates attacks
			\end{itemize}
		\end{block}
	\end{frame}
	
	% ================ SECTION 4: DOMAIN 4 - IT AND SECURITY ================
	
	\section{Domain 4: Enterprise Architecture}
	
	\begin{frame}{Q60: Technology Architecture}
		\begin{block}{Question}
			Networking devices, storage, and software are components of:
		\end{block}
		\begin{enumerate}[(A)]
			\item Business architecture
			\item Technology architecture
			\item Data architecture
			\item Application architecture
		\end{enumerate}
	\end{frame}
	
	\begin{frame}{Q60: Explanation}
		\begin{block}{Correct Answer: B}
			Networking devices, storage, and software are components of technology architecture.
		\end{block}
		\begin{block}{Option Analysis}
			\begin{itemize}
				\item \textbf{A - Incorrect:} Business architecture defines processes
				\item \textbf{B - Correct:} Technology architecture covers infrastructure
				\item \textbf{C - Incorrect:} Data architecture covers data management
				\item \textbf{D - Incorrect:} Application architecture covers software solutions
			\end{itemize}
		\end{block}
	\end{frame}
	
	\begin{frame}{Q61: CMM Defined Process}
		\begin{block}{Question}
			Which is NOT a property of the Defined process in CMM?
		\end{block}
		\begin{enumerate}[(A)]
			\item Well characterized and understood
			\item Defined at organizational level
			\item Proactive
			\item Unpredictable
		\end{enumerate}
	\end{frame}
	
	\begin{frame}{Q61: Explanation}
		\begin{block}{Correct Answer: D}
			Unpredictable is not a property of Defined process in CMM.
		\end{block}
		\begin{block}{Option Analysis}
			\begin{itemize}
				\item \textbf{A - True:} Well characterized
				\item \textbf{B - True:} Defined at organizational level
				\item \textbf{C - True:} Proactive
				\item \textbf{D - Not True:} Defined processes are predictable
			\end{itemize}
		\end{block}
	\end{frame}
	
	\begin{frame}{Q61: CMM Defined Process} \begin{block}{Question} Which of the following is NOT a characteristic of a Level 3 (Defined) process in the Capability Maturity Model (CMM)? \end{block} \begin{enumerate}[(A)] \item Processes are standardized across the organization \item Processes are documented and communicated \item Processes are primarily dependent on individual effort \item Processes are well understood by personnel \end{enumerate} 
	
	\begin{block}{Memory Aid} Level 1 = Ad Hoc Level 2 = Repeatable Level 3 = Defined Level 4 = Measured Level 5 = Optimized \end{block} 
	
	\end{frame} \begin{frame}{Q61: Explanation} \footnotesize \begin{block}{Correct Answer: C} A Level 3 (Defined) process is organizationally standardized, documented, and understood. Dependence on individual effort is characteristic of lower maturity levels. \end{block} \begin{block}{Option Analysis} \textbf{A - Incorrect} Standardization across the organization is a defining characteristic of Level 3 maturity. \vspace{0.15cm} \textbf{B - Incorrect} Defined processes are formally documented and communicated throughout the organization. \vspace{0.15cm} \textbf{C - Correct} Heavy reliance on individual effort indicates an ad hoc or immature process rather than a defined process. \vspace{0.15cm} \textbf{D - Incorrect} Personnel should understand and consistently follow defined processes. \end{block} \end{frame}
	
	
	\begin{frame}{Q62: TCP/IP Model Layers}
		\begin{block}{Question}
			Which is Layer 3 of the TCP/IP model?
		\end{block}
		\begin{enumerate}[(A)]
			\item Physical
			\item Data link
			\item Network
			\item Application
		\end{enumerate}
	\end{frame}
	
	\begin{frame}{Q62: Explanation}
		\begin{block}{Correct Answer: C}
			Network layer is Layer 3 of the TCP/IP model.
		\end{block}
		\begin{block}{Option Analysis}
			\begin{itemize}
				\item \textbf{A - Incorrect:} Physical is Layer 1
				\item \textbf{B - Incorrect:} Data link is Layer 2
				\item \textbf{C - Correct:} Network is Layer 3
				\item \textbf{D - Incorrect:} Application is Layer 5
			\end{itemize}
		\end{block}
	\end{frame}
	
	\begin{frame}{Q63: OSI Model Session Layer}
		\begin{block}{Question}
			Which is Layer 5 of the OSI model?
		\end{block}
		\begin{enumerate}[(A)]
			\item Physical
			\item Application
			\item Session
			\item Transport
		\end{enumerate}
	\end{frame}
	
	\begin{frame}{Q63: Explanation}
		\begin{block}{Correct Answer: C}
			Session layer is Layer 5 of the OSI model.
		\end{block}
		\begin{block}{Option Analysis}
			\begin{itemize}
				\item \textbf{A - Incorrect:} Physical is Layer 1
				\item \textbf{B - Incorrect:} Application is Layer 7
				\item \textbf{C - Correct:} Session is Layer 5
				\item \textbf{D - Incorrect:} Transport is Layer 4
			\end{itemize}
		\end{block}
	\end{frame}
	
	\begin{frame}{Q64: VPN Protocol}
		\begin{block}{Question}
			Which protocol is used to implement a VPN?
		\end{block}
		\begin{enumerate}[(A)]
			\item DNS
			\item IPSec
			\item SSL
			\item SSO
		\end{enumerate}
	\end{frame}
	
	\begin{frame}{Q64: Explanation}
		\begin{block}{Correct Answer: B}
			IPSec protocol is commonly used to implement VPNs.
		\end{block}
		\begin{block}{Option Analysis}
			\begin{itemize}
				\item \textbf{A - Incorrect:} DNS is domain name service
				\item \textbf{B - Correct:} IPSec enables VPN
				\item \textbf{C - Incorrect:} SSL is for web encryption
				\item \textbf{D - Incorrect:} SSO is single sign-on
			\end{itemize}
		\end{block}
	\end{frame}
	
	\begin{frame}{Q65: IaaS}
		\begin{block}{Question}
			Which is an example of on-demand compute, network, and storage services?
		\end{block}
		\begin{enumerate}[(A)]
			\item PaaS
			\item SaaS
			\item IaaS
			\item NaaS
		\end{enumerate}
	\end{frame}
	
	\begin{frame}{Q65: Explanation}
		\begin{block}{Correct Answer: C}
			IaaS provides on-demand compute, network, and storage services.
		\end{block}
		\begin{block}{Option Analysis}
			\begin{itemize}
				\item \textbf{A - Incorrect:} PaaS provides platform
				\item \textbf{B - Incorrect:} SaaS provides software
				\item \textbf{C - Correct:} IaaS provides infrastructure
				\item \textbf{D - Incorrect:} NaaS provides network services
			\end{itemize}
		\end{block}
	\end{frame}
	
	\begin{frame}{Q66: SaaS Example}
		\begin{block}{Question}
			An online platform for video games is an example of:
		\end{block}
		\begin{enumerate}[(A)]
			\item PaaS
			\item SaaS
			\item IaaS
			\item NaaS
		\end{enumerate}
	\end{frame}
	
	\begin{frame}{Q66: Explanation}
		\begin{block}{Correct Answer: B}
			Online video game platform that lets users play without installing software is SaaS.
		\end{block}
		\begin{block}{Option Analysis}
			\begin{itemize}
				\item \textbf{A - Incorrect:} PaaS provides development platform
				\item \textbf{B - Correct:} SaaS provides ready-to-use software
				\item \textbf{C - Incorrect:} IaaS provides infrastructure
				\item \textbf{D - Incorrect:} NaaS provides network
			\end{itemize}
		\end{block}
	\end{frame}
	
	\begin{frame}{Q67: Private Cloud}
		\begin{block}{Question}
			In which deployment model are underlying cloud services NOT shared with other customers?
		\end{block}
		\begin{enumerate}[(A)]
			\item Public
			\item Private
			\item Community
			\item Hybrid
		\end{enumerate}
	\end{frame}
	
	\begin{frame}{Q67: Explanation}
		\begin{block}{Correct Answer: B}
			Private cloud resources are dedicated to a single organization.
		\end{block}
		\begin{block}{Option Analysis}
			\begin{itemize}
				\item \textbf{A - Incorrect:} Public is shared
				\item \textbf{B - Correct:} Private is dedicated
				\item \textbf{C - Incorrect:} Community is shared with similar orgs
				\item \textbf{D - Incorrect:} Hybrid combines models
			\end{itemize}
		\end{block}
	\end{frame}
	
	\begin{frame}{Q68: Community Cloud}
		\begin{block}{Question}
			In which deployment model are underlying cloud services shared with similar customers?
		\end{block}
		\begin{enumerate}[(A)]
			\item Public
			\item Private
			\item Community
			\item Hybrid
		\end{enumerate}
	\end{frame}
	
	\begin{frame}{Q68: Explanation}
		\begin{block}{Correct Answer: C}
			Community cloud lets similar organizations share resources.
		\end{block}
		\begin{block}{Option Analysis}
			\begin{itemize}
				\item \textbf{A - Incorrect:} Public is open to all
				\item \textbf{B - Incorrect:} Private is dedicated
				\item \textbf{C - Correct:} Community is for similar orgs
				\item \textbf{D - Incorrect:} Hybrid combines models
			\end{itemize}
		\end{block}
	\end{frame}
	
	\section{Domain 4: Data Life Cycle}
	
	\begin{frame}{Q69: BIA Input for BCP}
		\begin{block}{Question}
			Which acts as an input for the BCP?
		\end{block}
		\begin{enumerate}[(A)]
			\item Recovery objectives
			\item DR test
			\item BIA
			\item BC test
		\end{enumerate}
	\end{frame}
	
	\begin{frame}{Q69: Explanation}
		\begin{block}{Correct Answer: C}
			BIA is the input to creating a BCP.
		\end{block}
		\begin{block}{Option Analysis}
			\begin{itemize}
				\item \textbf{A - Incorrect:} Recovery objectives are output of BIA
				\item \textbf{B - Incorrect:} DR test validates recovery
				\item \textbf{C - Correct:} BIA inputs BCP
				\item \textbf{D - Incorrect:} BC test validates BCP
			\end{itemize}
		\end{block}
	\end{frame}
	
	
	\begin{frame}{Q69: Business Continuity Planning} \begin{block}{Question} Which of the following is a primary input used when developing a Business Continuity Plan (BCP)? \end{block} \begin{enumerate}[(A)] \item Recovery Time Objective (RTO) \item Disaster Recovery Test Results \item Business Impact Analysis (BIA) \item Business Continuity Test Results \end{enumerate} \begin{block}{Memory Aid} BIA $\rightarrow$ BCP BIA outputs RTO, RPO, MTD BCP is built using BIA results BC/DR Tests validate the completed plan \end{block} \end{frame} \begin{frame}{Q69: Explanation} \footnotesize \begin{block}{Correct Answer: C - Business Impact Analysis (BIA)} The BIA identifies critical business processes, dependencies, recovery priorities, RTOs, RPOs, and business impacts. These findings provide the foundation for developing a Business Continuity Plan. \end{block} \begin{block}{Option Analysis} \textbf{A - Incorrect: Recovery Time Objective (RTO)} RTO is not an input to the BIA. Instead, it is one of the key outputs produced by the BIA and later incorporated into the BCP. \vspace{0.15cm} \textbf{B - Incorrect: Disaster Recovery Test Results} DR test results validate the effectiveness of recovery procedures after plans have already been developed. They improve and refine the BCP but are not primary inputs to its creation. \vspace{0.15cm} \textbf{C - Correct: Business Impact Analysis (BIA)} The BIA identifies critical processes, business impacts, dependencies, recovery requirements, and recovery objectives. These results are used to develop continuity and recovery strategies. \vspace{0.15cm} \textbf{D - Incorrect: Business Continuity Test Results} BC test results evaluate whether the BCP works as intended. Testing occurs after plan development and therefore cannot be considered a primary input. \end{block}  \end{frame}
	
	\begin{frame}{Q70: RPO}
		\begin{block}{Question}
			Risk manager confirms application can't lose more than four hours of data. Which objective is this?
		\end{block}
		\begin{enumerate}[(A)]
			\item MTD
			\item RPO
			\item RTO
			\item BIA
		\end{enumerate}
	\end{frame}
	
	\begin{frame}{Q70: Explanation}
		\begin{block}{Correct Answer: B}
			RPO is the maximum acceptable data loss - in this case, 4 hours.
		\end{block}
		\begin{block}{Option Analysis}
			\begin{itemize}
				\item \textbf{A - Incorrect:} MTD is total downtime
				\item \textbf{B - Correct:} RPO measures data loss tolerance
				\item \textbf{C - Incorrect:} RTO measures downtime tolerance
				\item \textbf{D - Incorrect:} BIA is the assessment process
			\end{itemize}
		\end{block}
	\end{frame}
	
	
	\begin{frame}{Q70: RPO} \begin{block}{Question} A risk manager confirms a critical application cannot lose more than four hours of data following a disruption. Which recovery objective does this requirement represent? \end{block} \begin{enumerate}[(A)] \item Maximum Tolerable Downtime (MTD) \item Recovery Point Objective (RPO) \item Recovery Time Objective (RTO) \item Business Impact Analysis (BIA) \end{enumerate} \begin{block}{Memory Aid} RPO = Data Loss ("How much data can we lose?") RTO = Downtime ("How long can systems be down?") MTD = Survival Limit ("How long can the business survive?") \end{block} \end{frame} \begin{frame}{Q70: Explanation} \footnotesize \begin{block}{Correct Answer: B - Recovery Point Objective (RPO)} RPO represents the maximum acceptable amount of data loss measured in time. If no more than four hours of data can be lost, the RPO is four hours. \end{block} \begin{block}{Option Analysis} \textbf{A - Incorrect: Maximum Tolerable Downtime (MTD)} MTD measures how long a business process can remain unavailable before unacceptable business damage occurs. It focuses on business survival, not data loss. \vspace{0.15cm} \textbf{B - Correct: Recovery Point Objective (RPO)} RPO measures allowable data loss. A four-hour RPO means recovery may require restoring data up to four hours before the disruption. \vspace{0.15cm} \textbf{C - Incorrect: Recovery Time Objective (RTO)} RTO measures how quickly a system must be restored after an outage. It focuses on downtime rather than lost data. \vspace{0.15cm} \textbf{D - Incorrect: Business Impact Analysis (BIA)} BIA is the process used to identify business impacts and establish recovery requirements such as RTO, RPO, and MTD. It is not itself a recovery objective. \end{block}  \end{frame}
	
	\begin{frame}{Q71: RTO and MTD}
		\begin{block}{Question}
			An application can tolerate four hours of downtime but must be recovered within six hours. Which metrics are correct?
		\end{block}
		\begin{enumerate}[(A)]
			\item RPO - 4h, RTO - 6h
			\item MTD - 4h, RTO - 6h
			\item MTD - 6h, RPO - 4h
			\item RTO - 4h, MTD - 6h
		\end{enumerate}
	\end{frame}
	
	\begin{frame}{Q71: Explanation}
		\begin{block}{Correct Answer: D}
			RTO is maximum acceptable downtime (4h). MTD is maximum total downtime before material impact (6h).
		\end{block}
		\begin{block}{Option Analysis}
			\begin{itemize}
				\item \textbf{A - Incorrect:} RPO is data loss, not downtime
				\item \textbf{B - Incorrect:} MTD should be 6h, not 4h
				\item \textbf{C - Incorrect:} MTD is 6h not 4h
				\item \textbf{D - Correct:} RTO=4h (downtime tolerance), MTD=6h (total before impact)
			\end{itemize}
		\end{block}
	\end{frame}
	
	
	\begin{frame}{Q71: RTO vs MTD} \begin{block}{Question} A Business Impact Analysis determines that a critical application must be recovered within four hours of an outage. The business can tolerate no more than six hours of total disruption before severe operational impact occurs. Which of the following is correct? \end{block} \begin{enumerate}[(A)] \item RPO = 4 hours, RTO = 6 hours \item MTD = 4 hours, RTO = 6 hours \item MTD = 6 hours, RPO = 4 hours \item RTO = 4 hours, MTD = 6 hours \end{enumerate} \begin{block}{Memory Aid} RPO = Data Loss RTO = Recovery Target MTD = Business Survival Limit Rule: RTO should always be less than or equal to MTD. \end{block} \end{frame} \begin{frame}{Q71: Explanation} \footnotesize \begin{block}{Correct Answer: D} RTO is the target recovery time for the system, while MTD is the maximum amount of downtime the business can tolerate before unacceptable consequences occur. \end{block} \begin{block}{Option Analysis} \textbf{A - Incorrect} RPO measures acceptable data loss, not downtime. The scenario provides no information regarding data loss. \vspace{0.15cm} \textbf{B - Incorrect} The application must be recovered within four hours, making the RTO four hours. The business can tolerate six hours of disruption, making the MTD six hours. \vspace{0.15cm} \textbf{C - Incorrect} MTD equals six hours, but there is no information concerning data loss, so RPO cannot be determined. \vspace{0.15cm} \textbf{D - Correct} RTO = 4 hours because recovery must occur within four hours. MTD = 6 hours because the business cannot tolerate disruption beyond six hours. \end{block}  \end{frame}
	
	\begin{frame}{Q72: Data Classification}
		\begin{block}{Question}
			Healthcare records and driver's licenses - most appropriate classification?
		\end{block}
		\begin{enumerate}[(A)]
			\item Sensitive (PII and PHI)
			\item Public
			\item PII
			\item PHI
		\end{enumerate}
	\end{frame}
	
	\begin{frame}{Q72: Explanation}
		\begin{block}{Correct Answer: A}
			Contains both PII and PHI, so classification should be Sensitive.
		\end{block}
		\begin{block}{Option Analysis}
			\begin{itemize}
				\item \textbf{A - Correct:} Sensitive covers both PII and PHI
				\item \textbf{B - Incorrect:} Public is for non-sensitive data
				\item \textbf{C - Incorrect:} PII alone doesn't cover health records
				\item \textbf{D - Incorrect:} PHI alone doesn't cover driver's licenses
			\end{itemize}
		\end{block}
	\end{frame}
	

	\begin{frame}{Q74: Data Labeling Purpose}
		\begin{block}{Question}
			Data labeling is important for which of the following?
		\end{block}
		\begin{enumerate}[(A)]
			\item Data archival
			\item Data destruction
			\item DLP
			\item All of the above
		\end{enumerate}
	\end{frame}
	
	\begin{frame}{Q74: Explanation}
		\begin{block}{Correct Answer: D}
			Data labeling is important for archival, destruction, and DLP.
		\end{block}
		\begin{block}{Option Analysis}
			\begin{itemize}
				\item \textbf{A - Incorrect:} Archival alone not complete
				\item \textbf{B - Incorrect:} Destruction alone not complete
				\item \textbf{C - Incorrect:} DLP alone not complete
				\item \textbf{D - Correct:} All functions benefit from labeling
			\end{itemize}
		\end{block}
	\end{frame}
	
	
	\begin{frame}{Q74: Data Classification} \begin{block}{Question} Why is data classification and labeling important from a risk management perspective? \end{block} \begin{enumerate}[(A)] \item It helps determine appropriate retention and archival requirements \item It helps determine appropriate disposal and destruction requirements \item It enables enforcement of data loss prevention (DLP) controls \item All of the above \end{enumerate}  \end{frame} \begin{frame}[shrink=10]{Q74: Explanation} \footnotesize \begin{block}{Correct Answer: D - All of the Above} Data labeling identifies the sensitivity and classification of information. Organizations rely on these labels throughout the information lifecycle to determine how data should be stored, protected, retained, shared, archived, and ultimately destroyed. \end{block} \begin{block}{Option Analysis} \textbf{A - Incorrect} Data classification assists with retention and archival decisions because different categories of information may be subject to different legal, regulatory, or business retention requirements. However, this is only one use of labeling. \vspace{0.15cm} \textbf{B - Incorrect} Data labels help determine destruction requirements. For example, confidential customer information may require secure deletion, whereas public information may not require special disposal procedures. However, this is not the only purpose of labeling. \vspace{0.15cm} \textbf{C - Incorrect} DLP technologies often rely on data classification labels to identify sensitive information and prevent unauthorized transmission, copying, or disclosure. However, labeling supports many other governance and security functions as well. \vspace{0.15cm} \textbf{D - Correct} Data labeling supports the entire information lifecycle, including protection, DLP enforcement, retention, archival, legal holds, access control decisions, and secure destruction. \end{block}  \end{frame}
	
	\section{Domain 4: SDLC}
	

	
	\begin{frame}{Q76: User Training Phase}
		\begin{block}{Question}
			End user training is an integral part of which phase of the SDLC?
		\end{block}
		\begin{enumerate}[(A)]
			\item Development
			\item Disposal
			\item Initiation
			\item Implementation
		\end{enumerate}
	\end{frame}
	
	\begin{frame}{Q76: Explanation}
		\begin{block}{Correct Answer: D}
			End users are trained in the implementation phase.
		\end{block}
		\begin{block}{Option Analysis}
			\begin{itemize}
				\item \textbf{A - Incorrect:} Development is coding
				\item \textbf{B - Incorrect:} Disposal is retirement
				\item \textbf{C - Incorrect:} Initiation is planning
				\item \textbf{D - Correct:} Implementation includes training
			\end{itemize}
		\end{block}
	\end{frame}
	
	
	\begin{frame}{Q76: SDLC Training} \begin{block}{Question} A new customer relationship management (CRM) system is ready for deployment. Before go-live, employees attend instructor-led sessions and complete training on the new business processes. During which SDLC phase is this activity MOST likely occurring? \end{block} \begin{enumerate}[(A)] \item Development \item Disposal \item Initiation \item Implementation \end{enumerate} \begin{block}{Memory Aid} Initiation = Approve Development = Build Implementation = Deploy and Train Maintenance = Operate Disposal = Retire \end{block}  \end{frame} \begin{frame}{Q76: Explanation} \footnotesize \begin{block}{Correct Answer: D - Implementation} User training is typically performed during the implementation phase when the system is being deployed and prepared for operational use. \end{block} \begin{block}{Option Analysis} \textbf{A - Incorrect: Development} Development focuses on designing, building, configuring, and testing the solution. End users are generally not trained until the system is ready for deployment. \vspace{0.15cm} \textbf{B - Incorrect: Disposal} Disposal occurs at the end of the system's lifecycle and includes decommissioning systems, data archival, and secure destruction activities. \vspace{0.15cm} \textbf{C - Incorrect: Initiation} Initiation includes feasibility analysis, budgeting, business-case development, and project approval. User training is not normally conducted during this phase. \vspace{0.15cm} \textbf{D - Correct: Implementation} Implementation includes deployment, data migration, user acceptance activities, end-user training, and transition to production operations. \end{block} \end{frame}
	
	\begin{frame}{Q77: SDLC Final Stage}
		\begin{block}{Question}
			Which is the final stage of the SDLC?
		\end{block}
		\begin{enumerate}[(A)]
			\item Development
			\item Disposal
			\item Initiation
			\item Maintenance
		\end{enumerate}
	\end{frame}
	
	\begin{frame}{Q77: Explanation}
		\begin{block}{Correct Answer: B}
			Disposal is the final phase of the SDLC.
		\end{block}
		\begin{block}{Option Analysis}
			\begin{itemize}
				\item \textbf{A - Incorrect:} Development is early phase
				\item \textbf{B - Correct:} Disposal is final phase
				\item \textbf{C - Incorrect:} Initiation is first phase
				\item \textbf{D - Incorrect:} Maintenance is ongoing
			\end{itemize}
		\end{block}
	\end{frame}
	
	\begin{frame}{Q78: Project Risk vs SDLC Risk}
		\begin{block}{Question}
			Key difference between project risk and SDLC risk is:
		\end{block}
		\begin{enumerate}[(A)]
			\item Project risk relates to objectives, SDLC relates to time
			\item Project risk relates to objectives, SDLC relates to development
			\item Project risk relates to development, SDLC relates to objectives
			\item There is no difference
		\end{enumerate}
	\end{frame}
	
	\begin{frame}{Q78: Explanation}
		\begin{block}{Correct Answer: B}
			Project risk relates to project objectives; SDLC risk relates to development objectives.
		\end{block}
		\begin{block}{Option Analysis}
			\begin{itemize}
				\item \textbf{A - Incorrect:} SDLC risk is not just time
				\item \textbf{B - Correct:} Different focuses
				\item \textbf{C - Incorrect:} Reversed relationship
				\item \textbf{D - Incorrect:} They are different
			\end{itemize}
		\end{block}
	\end{frame}
	
	
	\begin{frame}{Q78: Project Risk vs. SDLC Risk} \begin{block}{Question} What is the PRIMARY difference between project risk and SDLC risk? \end{block} \begin{enumerate}[(A)] \item Project risk affects project objectives, while SDLC risk affects the system being developed. \item Project risk affects system security, while SDLC risk affects project budget. \item Project risk occurs only during implementation, while SDLC risk occurs only during development. \item There is no difference between project risk and SDLC risk. \end{enumerate} \end{frame} \begin{frame}{Q78: Explanation} \footnotesize \begin{block}{Correct Answer: A} Project risk focuses on the successful completion of the project, while SDLC risk focuses on risks associated with the system, application, or solution being developed. \end{block} \begin{block}{Option Analysis} \textbf{A - Correct} Project risks threaten project objectives such as schedule, budget, resources, or scope. SDLC risks relate to weaknesses in the resulting solution, such as security vulnerabilities, availability concerns, or design flaws. \vspace{0.15cm} \textbf{B - Incorrect} This reverses the concepts. Security risks are generally SDLC risks, while budget risks are typically project risks. \vspace{0.15cm} \textbf{C - Incorrect} Both project risks and SDLC risks may exist throughout multiple phases of the project lifecycle. \vspace{0.15cm} \textbf{D - Incorrect} While related, they have different objectives and require different management approaches. \end{block}  \end{frame}
	
	\begin{frame}{Q79: Certification Requirement}
		\begin{block}{Question}
			Which is a must to achieve certification?
		\end{block}
		\begin{enumerate}[(A)]
			\item A risk manager
			\item An internal audit team
			\item Policies
			\item An external audit firm
		\end{enumerate}
	\end{frame}
	
	\begin{frame}{Q79: Explanation}
		\begin{block}{Correct Answer: D}
			An external audit firm is required to certify a product.
		\end{block}
		\begin{block}{Option Analysis}
			\begin{itemize}
				\item \textbf{A - Incorrect:} Risk manager may not be required
				\item \textbf{B - Incorrect:} Internal audit cannot certify externally
				\item \textbf{C - Incorrect:} Policies alone don't certify
				\item \textbf{D - Correct:} External audit firm is required
			\end{itemize}
		\end{block}
	\end{frame}
	
	
	
	
	
	\section{Domain 4: Information Security}
	
	\begin{frame}{Q82: Encryption Keys}
		\begin{block}{Question}
			Which keys are required for asymmetric key encryption?
		\end{block}
		\begin{enumerate}[(A)]
			\item Only public key
			\item Only private key
			\item Any public or private key
			\item Public and private key pair
		\end{enumerate}
	\end{frame}
	
	\begin{frame}{Q82: Explanation}
		\begin{block}{Correct Answer: D}
			Asymmetric encryption requires both a public key and a private key.
		\end{block}
		\begin{block}{Option Analysis}
			\begin{itemize}
				\item \textbf{A - Incorrect:} Public key alone insufficient
				\item \textbf{B - Incorrect:} Private key alone insufficient
				\item \textbf{C - Incorrect:} Need both keys
				\item \textbf{D - Correct:} Both keys are required
			\end{itemize}
		\end{block}
	\end{frame}
	
	\begin{frame}{Q83: Symmetric Encryption}
		\begin{block}{Question}
			Which key is required for symmetric key encryption?
		\end{block}
		\begin{enumerate}[(A)]
			\item Public and private key
			\item A single key for encryption and decryption
			\item Only public key and hash
			\item Only private key and hash
		\end{enumerate}
	\end{frame}
	
	\begin{frame}{Q83: Explanation}
		\begin{block}{Correct Answer: B}
			Symmetric encryption uses a single key for both encryption and decryption.
		\end{block}
		\begin{block}{Option Analysis}
			\begin{itemize}
				\item \textbf{A - Incorrect:} That's asymmetric
				\item \textbf{B - Correct:} Single key for both
				\item \textbf{C - Incorrect:} Public key is asymmetric
				\item \textbf{D - Incorrect:} Private key is asymmetric
			\end{itemize}
		\end{block}
	\end{frame}
	
	\begin{frame}{Q84: Integrity Definition}
		\begin{block}{Question}
			Integrity of information means that:
		\end{block}
		\begin{enumerate}[(A)]
			\item Information is available to authorized users only
			\item Information is available at all times
			\item Information has not been tampered with
			\item All of these
		\end{enumerate}
	\end{frame}
	
	\begin{frame}{Q84: Explanation}
		\begin{block}{Correct Answer: C}
			Integrity means information has not been tampered with.
		\end{block}
		\begin{block}{Option Analysis}
			\begin{itemize}
				\item \textbf{A - Incorrect:} That's confidentiality
				\item \textbf{B - Incorrect:} That's availability
				\item \textbf{C - Correct:} Integrity is no tampering
				\item \textbf{D - Incorrect:} Only C is correct
			\end{itemize}
		\end{block}
	\end{frame}
	
	\begin{frame}{Q85: Confidentiality Definition}
		\begin{block}{Question}
			Confidentiality of information means that:
		\end{block}
		\begin{enumerate}[(A)]
			\item Information is available to authorized users only
			\item Information is available at all times
			\item Information has not been tampered with
			\item All of these
		\end{enumerate}
	\end{frame}
	
	\begin{frame}{Q85: Explanation}
		\begin{block}{Correct Answer: A}
			Confidentiality means information is available to authorized users only.
		\end{block}
		\begin{block}{Option Analysis}
			\begin{itemize}
				\item \textbf{A - Correct:} Confidentiality limits access
				\item \textbf{B - Incorrect:} That's availability
				\item \textbf{C - Incorrect:} That's integrity
				\item \textbf{D - Incorrect:} Only A is correct
			\end{itemize}
		\end{block}
	\end{frame}
	
	\begin{frame}{Q86: Availability Definition}
		\begin{block}{Question}
			Availability of information means that:
		\end{block}
		\begin{enumerate}[(A)]
			\item Information is available to authorized users only
			\item Information is available at all times to authorized users
			\item Information has not been tampered with
			\item All of these
		\end{enumerate}
	\end{frame}
	
	\begin{frame}{Q86: Explanation}
		\begin{block}{Correct Answer: B}
			Availability means information is accessible when needed by authorized users.
		\end{block}
		\begin{block}{Option Analysis}
			\begin{itemize}
				\item \textbf{A - Incorrect:} That's confidentiality
				\item \textbf{B - Correct:} Availability is timely access
				\item \textbf{C - Incorrect:} That's integrity
				\item \textbf{D - Incorrect:} Only B is correct
			\end{itemize}
		\end{block}
	\end{frame}
	
	\begin{frame}{Q87: Accountability}
		\begin{block}{Question}
			Which access control principle speaks to logging and monitoring of user activities?
		\end{block}
		\begin{enumerate}[(A)]
			\item Identification
			\item Authentication
			\item Authorization
			\item Accountability
		\end{enumerate}
	\end{frame}
	
	\begin{frame}{Q87: Explanation}
		\begin{block}{Correct Answer: D}
			Accountability deals with logging and monitoring of user activities.
		\end{block}
		\begin{block}{Option Analysis}
			\begin{itemize}
				\item \textbf{A - Incorrect:} Identification is unique ID
				\item \textbf{B - Incorrect:} Authentication is verification
				\item \textbf{C - Incorrect:} Authorization is permissions
				\item \textbf{D - Correct:} Accountability is logging/monitoring
			\end{itemize}
		\end{block}
	\end{frame}
	
	
	\begin{frame}{Q87: Access Control Principles} \begin{block}{Question} A security investigator needs to determine which employee accessed sensitive customer records, what actions were performed, and when those actions occurred. Which access control principle BEST supports this requirement? \end{block} \begin{enumerate}[(A)] \item Identification \item Authentication \item Authorization \item Accountability \end{enumerate}  \end{frame} \begin{frame}{Q87: Explanation} \footnotesize \begin{block}{Correct Answer: D - Accountability} Accountability ensures user actions can be traced to a specific individual through logging, monitoring, audit trails, and activity reviews. \end{block} \begin{block}{Option Analysis} \textbf{A - Incorrect: Identification} Identification occurs when a user claims an identity, such as entering a user ID or email address. It answers the question: "Who are you?" Identification alone does not prove the user is actually that person. \vspace{0.15cm} \textbf{B - Incorrect: Authentication} Authentication verifies the claimed identity through passwords, MFA, biometrics, or certificates. It answers the question: "Can you prove who you are?" Authentication does not determine what the user can do after login. \vspace{0.15cm} \textbf{C - Incorrect: Authorization} Authorization determines which resources and actions a user is permitted to access. It answers the question: "What are you allowed to do?" It governs permissions rather than monitoring activity. \vspace{0.15cm} \textbf{D - Correct: Accountability} Accountability requires logging and monitoring so actions can be traced back to a specific user. It answers the question: "Who performed the action?" This principle supports investigations, audits, and nonrepudiation. \end{block}  \end{frame}
	
	\begin{frame}{Q88: Identification}
		\begin{block}{Question}
			Which access control principle requires each user to have a separate unique name?
		\end{block}
		\begin{enumerate}[(A)]
			\item Identification
			\item Authentication
			\item Authorization
			\item Accountability
		\end{enumerate}
	\end{frame}
	
	\begin{frame}{Q88: Explanation}
		\begin{block}{Correct Answer: A}
			Identification requires all users to have a unique ID.
		\end{block}
		\begin{block}{Option Analysis}
			\begin{itemize}
				\item \textbf{A - Correct:} Identification requires unique ID
				\item \textbf{B - Incorrect:} Authentication verifies identity
				\item \textbf{C - Incorrect:} Authorization sets permissions
				\item \textbf{D - Incorrect:} Accountability logs activity
			\end{itemize}
		\end{block}
	\end{frame}
	
	\begin{frame}{Q89: Segregation of Duties}
		\begin{block}{Question}
			SOD principle ensures that:
		\end{block}
		\begin{enumerate}[(A)]
			\item All admin access to a single user
			\item Access to approve transactions is distributed
			\item Confidentiality is maintained
			\item Identification is in place
		\end{enumerate}
	\end{frame}
	
	\begin{frame}{Q89: Explanation}
		\begin{block}{Correct Answer: B}
			SOD ensures no single user has sole access to approve transactions.
		\end{block}
		\begin{block}{Option Analysis}
			\begin{itemize}
				\item \textbf{A - Incorrect:} SOD prevents single control
				\item \textbf{B - Correct:} SOD distributes authority
				\item \textbf{C - Incorrect:} Confidentiality is separate
				\item \textbf{D - Incorrect:} Identification is separate
			\end{itemize}
		\end{block}
	\end{frame}
	
	\begin{frame}{Q90: Hashing Characteristics}
		\begin{block}{Question}
			Which is true about hashing?
		\end{block}
		\begin{enumerate}[(A)]
			\item Provides unique output for each input
			\item Is non-reversible
			\item Helps maintain integrity
			\item All of these
		\end{enumerate}
	\end{frame}
	
	\begin{frame}{Q90: Explanation}
		\begin{block}{Correct Answer: D}
			All options are true for hashing - unique output, non-reversible, and integrity.
		\end{block}
		\begin{block}{Option Analysis}
			\begin{itemize}
				\item \textbf{A - True:} Unique output
				\item \textbf{B - True:} Non-reversible
				\item \textbf{C - True:} Integrity verification
				\item \textbf{D - Correct:} All are true
			\end{itemize}
		\end{block}
	\end{frame}
	
	
	
	\begin{frame}{Q90: Hashing} \begin{block}{Question} Which statement BEST describes the primary purpose of hashing? \end{block} \begin{enumerate}[(A)] \item To provide confidentiality of data \item To verify data integrity \item To provide authorization \item To facilitate key exchange \end{enumerate}  \begin{block}{Memory Aid} Encryption = Confidentiality Hashing = Integrity Digital Signatures = Integrity + Nonrepudiation Access Controls = Authorization \end{block} \end{frame} \begin{frame}{Q90: Explanation} \footnotesize \begin{block}{Correct Answer: B - Verify Data Integrity} Hashing generates a fixed-length value that can be used to determine whether data has been altered. If the hash changes, the data has changed. \end{block} \begin{block}{Option Analysis} \textbf{A - Incorrect:} Confidentiality is primarily provided by encryption. Hashing is not designed to allow recovery of the original data and therefore is not used for protecting information in transit or storage. \vspace{0.15cm} \textbf{B - Correct:} Hashing is commonly used to verify integrity. Matching hash values indicate that data has likely not been modified. \vspace{0.15cm} \textbf{C - Incorrect:} Authorization determines what actions a user is permitted to perform after authentication. \vspace{0.15cm} \textbf{D - Incorrect:} Key exchange is performed using cryptographic key-management mechanisms rather than hashing. \end{block} \end{frame}
	
	\begin{frame}{Q91: Digital Signatures}
		\begin{block}{Question}
			Which is NOT provided by digital signatures?
		\end{block}
		\begin{enumerate}[(A)]
			\item Integrity
			\item Non-repudiation
			\item Authenticity
			\item Confidentiality
		\end{enumerate}
	\end{frame}
	
	\begin{frame}{Q91: Explanation}
		\begin{block}{Correct Answer: D}
			Digital signatures do NOT provide confidentiality on their own.
		\end{block}
		\begin{block}{Option Analysis}
			\begin{itemize}
				\item \textbf{A - Incorrect:} Integrity is provided
				\item \textbf{B - Incorrect:} Non-repudiation is provided
				\item \textbf{C - Incorrect:} Authenticity is provided
				\item \textbf{D - Correct:} Confidentiality is not provided
			\end{itemize}
		\end{block}
	\end{frame}
	
	\begin{frame}{Q92: Encryption and Confidentiality}
		\begin{block}{Question}
			Sender wants only receiver to decrypt data. Which should sender use?
		\end{block}
		\begin{enumerate}[(A)]
			\item Receiver's public key
			\item Sender's public key
			\item Receiver's private key
			\item Sender's private key
		\end{enumerate}
	\end{frame}
	
	\begin{frame}{Q92: Explanation}
		\begin{block}{Correct Answer: A}
			Use receiver's public key to encrypt so only receiver can decrypt with their private key.
		\end{block}
		\begin{block}{Option Analysis}
			\begin{itemize}
				\item \textbf{A - Correct:} Receiver's public key ensures only receiver decrypts
				\item \textbf{B - Incorrect:} Anyone with sender's public key
				\item \textbf{C - Incorrect:} Receiver's private key is not shared
				\item \textbf{D - Incorrect:} Sender's private key would allow anyone
			\end{itemize}
		\end{block}
	\end{frame}
	
	\begin{frame}{Q93: Signed and Encrypted}
		\begin{block}{Question}
			Which cryptographic mechanism provides confidentiality, integrity, and non-repudiation?
		\end{block}
		\begin{enumerate}[(A)]
			\item Signed, not encrypted
			\item Encrypted, not signed
			\item Signed, encrypted
			\item Signed, hashed
		\end{enumerate}
	\end{frame}
	
	\begin{frame}{Q93: Explanation}
		\begin{block}{Correct Answer: C}
			Signed and encrypted provides confidentiality, integrity, and non-repudiation.
		\end{block}
		\begin{block}{Option Analysis}
			\begin{itemize}
				\item \textbf{A - Incorrect:} No confidentiality
				\item \textbf{B - Incorrect:} No integrity/non-repudiation
				\item \textbf{C - Correct:} All three provided
				\item \textbf{D - Incorrect:} No confidentiality
			\end{itemize}
		\end{block}
	\end{frame}
	
	
	\begin{frame}{Q93: Information Protection} \begin{block}{Question} A customer submits sensitive information through an online portal. Management requires that unauthorized parties cannot read the data, that any modification can be detected, and that the sender cannot later deny submitting the information. Which cryptographic approach BEST satisfies these requirements? \end{block} \begin{enumerate}[(A)] \item Digitally signed but not encrypted \item Encrypted but not digitally signed \item Digitally signed and encrypted \item Digitally signed and hashed \end{enumerate} \begin{block}{Memory Aid} Encryption = Confidentiality Hashing = Integrity Digital Signature = Integrity + Authentication + Nonrepudiation Signed + Encrypted = All of the Above \end{block}  \end{frame} \begin{frame}[shrink=10]{Q93: Explanation} \footnotesize \begin{block}{Correct Answer: C - Digitally Signed and Encrypted} Encryption provides confidentiality, while a digital signature provides integrity, authenticity, and nonrepudiation. Together, they satisfy all requirements in the scenario. \end{block} \begin{block}{Option Analysis} \textbf{A - Incorrect: Digitally Signed but Not Encrypted} A digital signature helps verify the sender's identity and confirms that the message has not been altered. However, anyone who intercepts the message can still read its contents. Confidentiality is missing. \vspace{0.15cm} \textbf{B - Incorrect: Encrypted but Not Digitally Signed} Encryption protects confidentiality by preventing unauthorized disclosure. However, it does not provide nonrepudiation and does not prove who originated the message. \vspace{0.15cm} \textbf{C - Correct: Digitally Signed and Encrypted} Encryption protects confidentiality. The digital signature provides integrity, authenticity, and nonrepudiation. This combination delivers all security objectives listed in the question. \vspace{0.15cm} \textbf{D - Incorrect: Digitally Signed and Hashed} Digital signatures rely on hashing internally and provide integrity and nonrepudiation. However, hashing does not provide confidentiality because the underlying information remains readable. \end{block} \end{frame}
	
	\begin{frame}{Q94: PKI SPOF}
		\begin{block}{Question}
			Which could be a SPOF in the PKI ecosystem?
		\end{block}
		\begin{enumerate}[(A)]
			\item Certificate
			\item CA
			\item Requester
			\item Browser
		\end{enumerate}
	\end{frame}
	
	\begin{frame}{Q94: Explanation}
		\begin{block}{Correct Answer: B}
			CA could be a SPOF if its private key is compromised.
		\end{block}
		\begin{block}{Option Analysis}
			\begin{itemize}
				\item \textbf{A - Incorrect:} Individual certificate isn't SPOF
				\item \textbf{B - Correct:} CA compromise affects all
				\item \textbf{C - Incorrect:} Requester isn't SPOF
				\item \textbf{D - Incorrect:} Browser isn't SPOF
			\end{itemize}
		\end{block}
	\end{frame}
	
	
	\begin{frame}{Q94: PKI Trust Infrastructure} \begin{block}{Question} An organization discovers that the private key of its root Certificate Authority (CA) has been compromised. Security leadership determines that certificates issued throughout the environment can no longer be trusted. Which PKI component represents the MOST significant single point of failure (SPOF) in this scenario? \end{block} \begin{enumerate}[(A)] \item Certificate \item Certificate Authority (CA) \item Certificate Requester \item Browser \end{enumerate} \begin{block}{CRISC Memory Aid} Certificate = Identity CA = Trust Browser = Consumes Trust If the CA fails, trust in the PKI can fail. \end{block}  \end{frame} \begin{frame}[shrink=10]{Q94: Explanation} \footnotesize \begin{block}{Correct Answer: B - Certificate Authority (CA)} The CA serves as the trust anchor of the PKI ecosystem. If the CA or its private key is compromised, confidence in certificates issued by that CA may be lost, affecting a large portion of the trust infrastructure. \end{block} \begin{block}{Option Analysis} \textbf{A - Incorrect: Certificate} An individual certificate affects only the specific user, server, or application to which it was issued. Compromise of a single certificate is significant but does not generally undermine the entire PKI environment. \vspace{0.15cm} \textbf{B - Correct: Certificate Authority (CA)} The CA validates identities and issues certificates trusted throughout the environment. Because trust originates from the CA, compromise of the CA can invalidate confidence in numerous certificates and disrupt authentication, encryption, and digital-signature processes. \vspace{0.15cm} \textbf{C - Incorrect: Certificate Requester} The requester initiates the certificate issuance process but does not establish trust for other entities. A compromised requester affects a limited set of certificates rather than the entire PKI hierarchy. \vspace{0.15cm} \textbf{D - Incorrect: Browser} Browsers rely on trusted certificates but are not the source of trust within the PKI ecosystem. A browser issue affects individual users, whereas CA compromise can affect all parties relying on the PKI. \end{block} \end{frame}
	
	\begin{frame}{Q95: Security Awareness Primary Goal}
		\begin{block}{Question}
			Most important aspect of security awareness training is:
		\end{block}
		\begin{enumerate}[(A)]
			\item Train on latest threats
			\item Show compliance with audits
			\item Create security-conscious culture
			\item Upskill employees
		\end{enumerate}
	\end{frame}
	
	\begin{frame}{Q95: Explanation}
		\begin{block}{Correct Answer: C}
			Most important goal is to create a security-conscious culture.
		\end{block}
		\begin{block}{Option Analysis}
			\begin{itemize}
				\item \textbf{A - Incorrect:} Latest threats change
				\item \textbf{B - Incorrect:} Compliance is secondary
				\item \textbf{C - Correct:} Culture is primary
				\item \textbf{D - Incorrect:} Upskilling is secondary
			\end{itemize}
		\end{block}
	\end{frame}

	\begin{frame}{Q97: Data Security vs Privacy}
		\begin{block}{Question}
			Which statement is true about data security and data privacy?
		\end{block}
		\begin{enumerate}[(A)]
			\item They are the same thing
			\item Security is protecting data; privacy is protecting personal information
			\item Privacy is protecting data; security is protecting personal information
			\item Not important for organizations
		\end{enumerate}
	\end{frame}
	
	\begin{frame}{Q97: Explanation}
		\begin{block}{Correct Answer: B}
			Data security protects data from unauthorized access; data privacy protects personal information from misuse.
		\end{block}
		\begin{block}{Option Analysis}
			\begin{itemize}
				\item \textbf{A - Incorrect:} They are different
				\item \textbf{B - Correct:} Security vs privacy distinction
				\item \textbf{C - Incorrect:} Reversed
				\item \textbf{D - Incorrect:} Both are important
			\end{itemize}
		\end{block}
	\end{frame}
	
	
	\begin{frame}{Q97: Data Security vs Privacy} \begin{block}{Question} A bank encrypts customer records, restricts access based on job responsibilities, and monitors unauthorized access attempts. However, regulators determine the bank is collecting more personal information than necessary for its stated business purpose. What does this situation BEST demonstrate? \end{block} \begin{enumerate}[(A)] \item Data security and data privacy are the same concept \item Strong data security does not necessarily guarantee data privacy \item Data privacy is more important than data security \item Data security applies only to personal information \end{enumerate} 
	\begin{block}{Memory Aid} Security = Protect the Data Privacy = Use the Data Properly You can have Security without Privacy. You cannot have Privacy without Security. \end{block} 
	\end{frame} \begin{frame}{Q97: Explanation} \footnotesize \begin{block}{Correct Answer: B - Strong Data Security Does Not Necessarily Guarantee Data Privacy} Security and privacy are related but distinct concepts. An organization may effectively protect information from unauthorized access while still violating privacy principles through excessive collection, inappropriate use, or improper sharing of personal information. \end{block} \begin{block}{Option Analysis} \textbf{A - Incorrect} Data security focuses on protecting information through controls such as encryption, access management, monitoring, and incident response. Data privacy focuses on the appropriate collection, use, sharing, retention, and disposal of personal information. \vspace{0.15cm} \textbf{B - Correct} The organization has implemented strong security controls, yet it may still have a privacy issue because it is collecting more personal information than necessary. This illustrates that security supports privacy but does not guarantee it. \vspace{0.15cm} \textbf{C - Incorrect} Both security and privacy are important and complementary. Effective governance programs must address both. \vspace{0.15cm} \textbf{D - Incorrect} Security protects all forms of information, including intellectual property, financial records, operational data, and personal information. Privacy primarily focuses on personal information. \end{block} \end{frame}
	
	\begin{frame}{Q98: Privacy Compliance}
		\begin{block}{Question}
			Company collecting customer data for marketing. Which measure ensures compliance with privacy regulations?
		\end{block}
		\begin{enumerate}[(A)]
			\item Use encryption
			\item Obtain consent
			\item Implement firewalls
			\item Train employees
		\end{enumerate}
	\end{frame}
	
	\begin{frame}{Q98: Explanation}
		\begin{block}{Correct Answer: B}
			Obtaining consent is critical for compliance with GDPR, CCPA, and other privacy regulations.
		\end{block}
		\begin{block}{Option Analysis}
			\begin{itemize}
				\item \textbf{A - Incorrect:} Encryption is security, not privacy
				\item \textbf{B - Correct:} Consent is required
				\item \textbf{C - Incorrect:} Firewalls are security
				\item \textbf{D - Incorrect:} Training is important but not primary
			\end{itemize}
		\end{block}
	\end{frame}
	
	
	\begin{frame}{Q98: Privacy Compliance} \begin{block}{Question} An organization plans to use customer information collected through its website to support targeted marketing campaigns. Before the information is used, management wants to reduce the risk of privacy compliance violations. Which action is MOST important? \end{block} \begin{enumerate}[(A)] \item Encrypt the customer database \item Obtain appropriate customer consent \item Deploy additional firewall controls \item Conduct security awareness training \end{enumerate}  \begin{block}{CRISC Memory Aid} Security asks: "How do we protect the data?" Privacy asks: "Do we have the right to collect and use the data?" Consent = Privacy Encryption/Firewalls = Security \end{block}  \end{frame} \begin{frame}{Q98: Explanation} \footnotesize \begin{block}{Correct Answer: B - Obtain Appropriate Customer Consent} Privacy regulations generally require organizations to have a lawful basis for collecting and processing personal information. Obtaining appropriate consent is often a key requirement when personal information is used for marketing purposes. \end{block} \begin{block}{Option Analysis} \textbf{A - Incorrect: Encrypt the Customer Database} Encryption helps protect the confidentiality of customer information and reduces security risk. However, strong security controls alone do not provide authorization to collect or use personal information for marketing purposes. \vspace{0.15cm} \textbf{B - Correct: Obtain Appropriate Customer Consent} Consent helps ensure that customers understand how their information will be used. From a privacy perspective, this is generally more important than technical controls when the issue involves the lawful use of personal information. \vspace{0.15cm} \textbf{C - Incorrect: Deploy Additional Firewall Controls} Firewalls help protect systems from unauthorized access and cyberattacks. They support security objectives but do not address whether the organization has permission to process personal data. \vspace{0.15cm} \textbf{D - Incorrect: Conduct Security Awareness Training} Training helps employees understand privacy and security responsibilities. While valuable, it does not by itself establish a legal basis for collecting or using customer information. \end{block} \end{frame}
	
	\begin{frame}{Q99: Vulnerability Management KPI}
		\begin{block}{Question}
			Which KPI is most appropriate to measure vulnerability management program effectiveness?
		\end{block}
		\begin{enumerate}[(A)]
			\item Number of vulnerabilities found and remediated
			\item Number of security incidents reported
			\item Percentage of employees trained
			\item Average incident response time
		\end{enumerate}
	\end{frame}
	
	\begin{frame}{Q99: Explanation}
		\begin{block}{Correct Answer: A}
			Number of vulnerabilities found and remediated measures vulnerability management effectiveness.
		\end{block}
		\begin{block}{Option Analysis}
			\begin{itemize}
				\item \textbf{A - Correct:} Direct measure of VM program
				\item \textbf{B - Incorrect:} Measures incident response
				\item \textbf{C - Incorrect:} Measures training program
				\item \textbf{D - Incorrect:} Measures IR effectiveness
			\end{itemize}
		\end{block}
	\end{frame}
	
	\begin{frame}{Q100: Security Program KPI}
		\begin{block}{Question}
			Which KPI is most useful to evaluate security program effectiveness over time?
		\end{block}
		\begin{enumerate}[(A)]
			\item Number of employees trained
			\item Number of security incidents reported
			\item Average detection and response time
			\item Percentage reduction in data breaches
		\end{enumerate}
	\end{frame}
	
	\begin{frame}{Q100: Explanation}
		\begin{block}{Correct Answer: D}
			Percentage reduction in data breaches directly measures program impact over time.
		\end{block}
		\begin{block}{Option Analysis}
			\begin{itemize}
				\item \textbf{A - Incorrect:} Training completion not a direct measure
				\item \textbf{B - Incorrect:} Incidents may change
				\item \textbf{C - Incorrect:} Response time is operational
				\item \textbf{D - Correct:} Direct measure of program success
			\end{itemize}
		\end{block}
	\end{frame}
	
	
	\begin{frame}[shrink=10]{Q100: Security Program Performance}
		\begin{block}{Question}
			Senior management wants to evaluate whether the organization's
			security program is becoming more effective at identifying and
			responding to security incidents.
			
			Which KPI would provide the MOST useful measurement?
		\end{block}
		
		\begin{enumerate}[(A)]
			\item Number of employees completing security awareness training
			\item Number of security incidents reported
			\item Average detection and response time
			\item Number of audit findings
		\end{enumerate}
	\end{frame}
	
	\begin{frame}[shrink=10]{Q100: Explanation}
		\footnotesize
		
		\begin{block}{Correct Answer: C - Average Detection and Response Time}
			Effective security programs should improve the organization's
			ability to detect and respond to threats quickly. Measuring
			detection and response time provides a direct indication of
			operational security effectiveness and can be tracked over time.
		\end{block}
		
		\begin{block}{Option Analysis}
			
			\textbf{A - Incorrect: Number of Employees Completing Security Awareness Training}
			
			Training completion measures participation in a security
			initiative rather than the effectiveness of the security
			program itself. Employees may complete training while
			security performance remains unchanged.
			

			
			\textbf{B - Incorrect: Number of Security Incidents Reported}
			
			A higher number of reported incidents may indicate either
			worsening security or improved detection and reporting.
			The metric alone does not reliably measure effectiveness.
			

			
			\textbf{C - Correct: Average Detection and Response Time}
			
			Faster detection and containment reduce organizational
			exposure and financial impact. Improvements in this KPI
			typically indicate stronger security operations.
			

			
			\textbf{D - Incorrect: Number of Audit Findings}
			
			Audit findings identify control weaknesses but may vary
			based on audit scope, timing, and methodology. They are
			not a direct measure of day-to-day security performance.
			
		\end{block}
		

		
	\end{frame}
	
	% ================ FINAL SLIDES ================
	
	\begin{frame}{CRISC Key Takeaways}
		\begin{block}{Essential Concepts}
			\begin{itemize}
				\item Governance provides oversight and direction from Board
				\item Risk identification is the first step - "You can't protect what you don't know"
				\item Risk Appetite = willingness to accept risk
				\item Risk Tolerance = acceptable variation from appetite
				\item Risk Capacity = maximum risk before existence threatened
				\item RPO is backward-looking; RTO is forward-looking
			\end{itemize}
		\end{block}
	\end{frame}
	
	\begin{frame}{CRISC Key Takeaways (continued)}
		\begin{block}{More Essential Concepts}
			\begin{itemize}
				\item 3LoD: First = Operations, Second = Oversight, Third = Audit
				\item Residual Risk = Inherent Risk - Controls
				\item Risk Response: Mitigate, Accept, Transfer, Avoid
				\item SMART: Specific, Measurable, Attainable, Relevant, Timely
				\item CIA: Confidentiality, Integrity, Availability
				\item Digital signatures: Integrity, Non-repudiation, Authenticity (NOT confidentiality)
			\end{itemize}
		\end{block}
	\end{frame}
	
	\begin{frame}[shrink=10]{Final Advice}
		\begin{block}{Exam Preparation}
			\centering
			\vspace{0.3cm}
			\textbf{Remember:}
			\begin{itemize}
				\item Understand the \textit{why} behind each concept
				\item Know control categories and their differences
				\item Understand recovery objectives (RPO, RTO, MTD)
				\item Know risk appetite vs tolerance vs capacity
				\item Practice scenario-based questions
			\end{itemize}
			\vspace{0.5cm}
			\textbf{Good Luck on Your CRISC Exam!}
		\end{block}
	\end{frame}
	
	
	
	
	\begin{frame}[shrink=10]{CRISC Memory Sheet: Risk Terms}
		\begin{block}{Risk Fundamentals}
			
			RA = Risk Assessment
			
			BIA = Business Impact Analysis
			
			ERM = Enterprise Risk Management
			
			KRIs = Key Risk Indicators
			
			KPIs = Key Performance Indicators
			
			KCIs = Key Control Indicators
			
			\vspace{0.15cm}
			
			Inherent Risk
			= Risk before controls.
			
			Residual Risk
			= Risk after controls.
			
			Control Gap
			= Missing or ineffective control.
			
			\vspace{0.15cm}
			
			Risk Appetite
			= Desired risk.
			
			Risk Tolerance
			= Acceptable variation.
			
			Risk Capacity
			= Maximum survivable risk.
			
		\end{block}
	\end{frame}
	
	\begin{frame}[shrink=10]{CRISC Memory Sheet: BIA}
		\begin{block}{Business Impact Analysis (BIA)}
			
			BIA Outputs:
			
			RTO = Recovery Time Objective
			
			RPO = Recovery Point Objective
			
			MTD = Maximum Tolerable Downtime
			
			MAO = Maximum Acceptable Outage
			
			\vspace{0.15cm}
			
			RPO = Data Loss
			
			RTO = Recovery Time
			
			MTD = Business Survival Limit
			
			\vspace{0.15cm}
			
			Rule:
			
			RTO $\leq$ MTD
			
			\vspace{0.15cm}
			
			BIA determines:
			
			Critical Processes
			
			Dependencies
			
			Business Impact
			
			Recovery Priorities
			
		\end{block}
	\end{frame}
	
	\begin{frame}[shrink=10]{CRISC Memory Sheet: Controls}
		\begin{block}{Control Types}
			
			Preventive = Stop
			
			Detective = Detect
			
			Corrective = Recover
			
			Deterrent = Discourage
			
			Compensating = Alternate
			
			Directive = Guide
			
			Recovery = Restore
			
			\vspace{0.15cm}
			
			MFA = Preventive
			
			Firewall = Preventive
			
			SIEM = Detective
			
			IDS = Detective
			
			Backup Restore = Corrective
			
			Security Banner = Deterrent
			
		\end{block}
	\end{frame}
	
	\begin{frame}[shrink=10]{CRISC Memory Sheet: Access Control}
		\begin{block}{IAAA Model}
			
			Identification
			
			= Who are you?
			
			\vspace{0.1cm}
			
			Authentication
			
			= Prove it.
			
			\vspace{0.1cm}
			
			Authorization
			
			= What can you do?
			
			\vspace{0.1cm}
			
			Accountability
			
			= What did you do?
			
			\vspace{0.15cm}
			
			Logging + Monitoring
			
			= Accountability
			
		\end{block}
	\end{frame}
	
	\begin{frame}[shrink=10]{CRISC Memory Sheet: Cryptography}
		\begin{block}{Crypto Fundamentals}
			
			PKI = Public Key Infrastructure
			
			CA = Certificate Authority
			
			CRL = Certificate Revocation List
			
			CSR = Certificate Signing Request
			
			\vspace{0.15cm}
			
			Encryption
			
			= Confidentiality
			
			Hashing
			
			= Integrity
			
			Digital Signature
			
			= Integrity + Authenticity + Nonrepudiation
			
			\vspace{0.15cm}
			
			Signed + Encrypted
			
			= Confidentiality + Integrity + Nonrepudiation
			
		\end{block}
	\end{frame}
	
	\begin{frame}[shrink=10]{CRISC Memory Sheet: Vendor Governance}
		\begin{block}{Third-Party Management}
			
			MSA = Master Service Agreement
			
			SLA = Service Level Agreement
			
			NDA = Non-Disclosure Agreement
			
			SOC = System and Organization Controls Report
			
			\vspace{0.15cm}
			
			MSA = Relationship
			
			SLA = Performance
			
			NDA = Confidentiality
			
			SOC = Assurance
			
			\vspace{0.15cm}
			
			Service Credits?
			
			= SLA
			
		\end{block}
	\end{frame}
	
	\begin{frame}[shrink=10]{CRISC Memory Sheet: SDLC}
		\begin{block}{System Development Life Cycle}
			
			Initiation
			
			= Budget + Approval
			
			\vspace{0.1cm}
			
			Development
			
			= Build + Test
			
			\vspace{0.1cm}
			
			Implementation
			
			= Deploy + Train
			
			\vspace{0.1cm}
			
			Maintenance
			
			= Operate + Patch
			
			\vspace{0.1cm}
			
			Disposal
			
			= Retire + Destroy
			
		\end{block}
	\end{frame}
	
	\begin{frame}[shrink=10]{CRISC Memory Sheet: CMM}
		\begin{block}{Capability Maturity Model}
			
			Level 1 = Initial
			
			(Ad Hoc)
			
			\vspace{0.1cm}
			
			Level 2 = Repeatable
			
			\vspace{0.1cm}
			
			Level 3 = Defined
			
			\vspace{0.1cm}
			
			Level 4 = Managed
			
			(Measured)
			
			\vspace{0.1cm}
			
			Level 5 = Optimized
			
			(Continuous Improvement)
			
		\end{block}
	\end{frame}
	
	\begin{frame}[shrink=10]{CRISC Memory Sheet: Privacy}
		\begin{block}{Privacy and Security}
			
			PII = Personally Identifiable Information
			
			PHI = Protected Health Information
			
			DLP = Data Loss Prevention
			
			\vspace{0.15cm}
			
			Security
			
			= Protect Data
			
			\vspace{0.1cm}
			
			Privacy
			
			= Proper Use of Data
			
			\vspace{0.1cm}
			
			Consent
			
			= Privacy
			
			\vspace{0.1cm}
			
			Encryption
			
			= Security
			
		\end{block}
	\end{frame}
	
	\begin{frame}[shrink=10]{Ultimate CRISC 30-Second Dump}
		\begin{block}{High-Yield Facts}
			
			BIA $\rightarrow$ RTO, RPO, MTD
			
			RA $\rightarrow$ Threats, Vulnerabilities
			
			KPI $\rightarrow$ Performance
			
			KRI $\rightarrow$ Risk
			
			KCI $\rightarrow$ Controls
			
			\vspace{0.15cm}
			
			Preventive = Stop
			
			Detective = Detect
			
			Corrective = Fix
			
			\vspace{0.15cm}
			
			Accept = Keep
			
			Mitigate = Control
			
			Transfer = Insurance
			
			Avoid = Stop Activity
			
			\vspace{0.15cm}
			
			Encryption = Confidentiality
			
			Hashing = Integrity
			
			Signature = Nonrepudiation
			
		\end{block}
	\end{frame}
	
\end{document}